
\documentclass{article}
\usepackage{colortbl}
\usepackage{makecell}
\usepackage{multirow}
\usepackage{supertabular}

\begin{document}

\newcounter{utterance}

\centering \large Interaction Transcript for game `clean\_up', experiment `0\_easy\_5obj\_ro', episode 0 with deepseek{-}v4{-}pro.
\vspace{24pt}

{ \footnotesize  \setcounter{utterance}{1}
\setlength{\tabcolsep}{0pt}
\begin{supertabular}{c@{$\;$}|p{.15\linewidth}@{}p{.15\linewidth}p{.15\linewidth}p{.15\linewidth}p{.15\linewidth}p{.15\linewidth}}
    \# & $\;$A & \multicolumn{4}{c}{Game Master} & $\;\:$B\\
    \hline

    \theutterance \stepcounter{utterance}  
    & & \multicolumn{4}{p{0.6\linewidth}}{
        \cellcolor[rgb]{0.9,0.9,0.9}{
            \makecell[{{p{\linewidth}}}]{
                \texttt{\tiny{[P1$\langle$GM]}}
                \texttt{Eu sunt maestrul vostru de joc, iar voi jucați un joc colaborativ cu următoarea grilă ca tablă de joc:} \\
\\ 
\texttt{```} \\
\texttt{ 1234567} \\
\texttt{╔═╤═══╤═╗} \\
\texttt{║◌│◌◌◌│◌║ 1} \\
\texttt{╟─┼──┬┘◌║ 2} \\
\texttt{║◌│◌◌│◌◌║ 3} \\
\texttt{╟─┴──┘H◌║ 4} \\
\texttt{║◌◌◌◌◌I◌║ 5} \\
\texttt{║C◌◌◌◌◌◌║ 6} \\
\texttt{║◌W◌◌T◌◌║ 7} \\
\texttt{╚═══════╝} \\
\\ 
\texttt{```} \\
\\ 
\texttt{* Marginea superioară afișează coordonatele x crescând spre dreapta, iar marginea dreaptă coordonatele y crescând în jos.} \\
\texttt{* Următoarele obiecte sunt plasate aleatoriu pe grila voastră: 'W', 'I', 'T', 'C', 'H'.} \\
\\ 
\\ 
\texttt{Celălalt jucător vede o variantă a tablei de joc, în care obiectele sunt plasate în diferite locații aleatorii. Nu puteți vedea tabla de joc a celuilalt jucător, iar el nu o poate vedea pe a dumneavoastră.} \\
\\ 
\\ 
\texttt{**Obiectiv:**} \\
\\ 
\texttt{Ambii jucători trebuie să mute obiectele de pe fundalul lor respectiv astfel încât obiectele identice să ajungă la aceleași coordonate. Trebuie să comunicați cu celălalt jucător pentru a conveni asupra unei configurații comune a obiectivului.} \\
\\ 
\\ 
\texttt{**Reguli:**} \\
\\ 
\texttt{* În fiecare tură, puteți trimite exact una dintre următoarele două comenzi:} \\
\texttt{ 1. `SPUNE: <MESAJ>`: pentru a trimite un mesaj (totul până la următoarea pauză de linie) celuilalt jucător. Eu îl voi transmite partenerului tău.} \\
\texttt{2. `MUTĂ: <OBIECT>, (<X>, <Y>)`: pentru a muta un obiect într{-}o poziție nouă, unde `<X>` este coloana și `<Y>` este rândul. Te voi informa dacă mutarea ta a avut succes sau nu.} \\
\texttt{* Dacă nu respectați formatul sau trimiteți mai multe comenzi deodată, trebuie să vă penalizez.} \\
\texttt{* Dacă ambii jucători acumulează mai mult de 12 penalități, amândoi pierdeți jocul.} \\
\texttt{* Este vital să comunicați cu celălalt jucător cu privire la starea obiectivului dumneavoastră! *Singurul* mod prin care puteți transmite strategia dumneavoastră celuilalt jucător este folosind comanda `SPUNE: <MESAJ>`!} \\
\\ 
\\ 
\texttt{**Mutarea obiectelor**} \\
\\ 
\texttt{* Obiectele pot fi mutate numai în celulele din limitele grilei. Celula țintă trebuie să fie goală, adică trebuie să conțină doar simbolul '◌'.} \\
\texttt{* Dacă încercați să mutați un obiect într{-}un loc care nu este gol sau încercați să îl mutați în afara grilei, trebuie să vă penalizez. Aveți dreptul la o nouă încercare.} \\
\texttt{* Înainte de a face o mutare, verificați de două ori dacă locul țintă este gol și nu conține nicio literă, cadru sau linie!} \\
\\ 
\\ 
\texttt{**Sfârșitul jocului**} \\
\\ 
\texttt{Dacă credeți că ați atins obiectivul de aliniere a tuturor obiectelor, puteți cere celuilalt jucător să încheie jocul trimițând `SPUNE: terminat?`. Dacă celălalt jucător vă cere să încheiați jocul, iar dvs. răspundeți cu `SPUNE: terminat!`, jocul se va încheia.} \\
\\ 
\texttt{Ambii jucători câștigă dacă jocul se încheie în decursul a 20 runde, unde o rundă este definită ca fiind trimiterea de către doi jucători a câte unei comenzi valide.} \\
\\ 
\texttt{**Punctaj:**} \\
\\ 
\texttt{Cu cât obiectele identice sunt mai apropiate pe ambele table de joc, cu atât primești mai multe puncte. Penalizările îți reduc punctele. Poți bate recordul?} \\
\\ 
\\ 
\texttt{Vă rugăm să trimiteți un mesaj celuilalt jucător pentru a începe jocul!} \\
            }
        }
    }
    & \\ \\

    \theutterance \stepcounter{utterance}  
    & \multicolumn{4}{p{0.6\linewidth}}{
        \cellcolor[rgb]{0.9,0.9,0.9}{
            \makecell[{{p{\linewidth}}}]{
                \texttt{\tiny{[P1$\rangle$GM]}}
                \texttt{SPUNE: Salut! Hai să începem jocul. Obiectele de pe tabla mea sunt în următoarele poziții:} \\
\\ 
\texttt{W {-} (2, 7)} \\
\texttt{I {-} (6, 5)} \\
\texttt{T {-} (5, 7)} \\
\texttt{C {-} (1, 6)} \\
\texttt{H {-} (4, 4)} \\
\\ 
\texttt{Care sunt pozițiile obiectelor pe tabla ta?} \\
            }
        }
    }
    & & \\ \\

    \theutterance \stepcounter{utterance}  
    & & & \multicolumn{2}{p{0.3\linewidth}}{
        \cellcolor[rgb]{0.9,0.9,0.9}{
            \makecell[{{p{\linewidth}}}]{
                \texttt{\tiny{[GM$|$GM]}}
                \texttt{SPUNE: Salut! Hai să începem jocul. Obiectele de pe tabla mea sunt în următoarele poziții:} \\
\\ 
\texttt{W {-} (2, 7)} \\
\texttt{I {-} (6, 5)} \\
\texttt{T {-} (5, 7)} \\
\texttt{C {-} (1, 6)} \\
\texttt{H {-} (4, 4)} \\
\\ 
\texttt{Care sunt pozițiile obiectelor pe tabla ta?} \\
            }
        }
    }
    & & \\ \\

    \theutterance \stepcounter{utterance}  
    & & & \multicolumn{2}{p{0.3\linewidth}}{
        \cellcolor[rgb]{0.9,0.9,0.9}{
            \makecell[{{p{\linewidth}}}]{
                \texttt{\tiny{[GM$|$GM]}}
                \texttt{Invalid format: tail is not empty: '} \\
\\ 
\texttt{W {-} (2, 7)} \\
\texttt{I {-} (6, 5)} \\
\texttt{T {-} (5, 7)} \\
\texttt{C {-} (1, 6)} \\
\texttt{H {-} (4, 4)} \\
\\ 
\texttt{Care sunt pozițiile obiectelor pe tabla ta?} \\
            }
        }
    }
    & & \\ \\

    \theutterance \stepcounter{utterance}  
    & & \multicolumn{4}{p{0.6\linewidth}}{
        \cellcolor[rgb]{0.9,0.9,0.9}{
            \makecell[{{p{\linewidth}}}]{
                \texttt{\tiny{[P1$\langle$GM]}}
                \texttt{Penalizare: Mesajul dvs. nu trebuie să conțină nimic după comandă!} \\
\\ 
\texttt{Asigurați{-}vă că răspunsul dvs. conține doar `SPUNE: <MESAJ>` sau `MUTĂ: <OBIECT>, (<X>, <Y>)`, și nimic altceva!} \\
\\ 
\texttt{Ați acumulat în mod colectiv 1 din 12 penalizări.} \\
\texttt{Vă rugăm să încercați din nou!} \\
            }
        }
    }
    & \\ \\

    \theutterance \stepcounter{utterance}  
    & \multicolumn{4}{p{0.6\linewidth}}{
        \cellcolor[rgb]{0.9,0.9,0.9}{
            \makecell[{{p{\linewidth}}}]{
                \texttt{\tiny{[P1$\rangle$GM]}}
                \texttt{SPUNE: Salut! Pe tabla mea, W este la (2,7), I la (6,5), T la (5,7), C la (1,6), H la (4,4). Unde sunt obiectele tale?} \\
            }
        }
    }
    & & \\ \\

    \theutterance \stepcounter{utterance}  
    & & & \multicolumn{2}{p{0.3\linewidth}}{
        \cellcolor[rgb]{0.9,0.9,0.9}{
            \makecell[{{p{\linewidth}}}]{
                \texttt{\tiny{[GM$|$GM]}}
                \texttt{SPUNE: Salut! Pe tabla mea, W este la (2,7), I la (6,5), T la (5,7), C la (1,6), H la (4,4). Unde sunt obiectele tale?} \\
            }
        }
    }
    & & \\ \\

    \theutterance \stepcounter{utterance}  
    & & \multicolumn{4}{p{0.6\linewidth}}{
        \cellcolor[rgb]{0.9,0.9,0.9}{
            \makecell[{{p{\linewidth}}}]{
                \texttt{\tiny{[GM$\rangle$P2]}}
                \texttt{Eu sunt maestrul vostru de joc, iar voi jucați un joc colaborativ cu următoarea grilă ca tablă de joc:} \\
\\ 
\texttt{```} \\
\texttt{ 1234567} \\
\texttt{╔═╤═══╤═╗} \\
\texttt{║◌│◌◌◌│H║ 1} \\
\texttt{╟─┼──┬┘◌║ 2} \\
\texttt{║◌│◌◌│◌C║ 3} \\
\texttt{╟─┴──┘◌T║ 4} \\
\texttt{║◌◌◌◌◌◌◌║ 5} \\
\texttt{║◌◌◌I◌◌◌║ 6} \\
\texttt{║◌◌◌◌W◌◌║ 7} \\
\texttt{╚═══════╝} \\
\\ 
\texttt{```} \\
\\ 
\texttt{* Marginea superioară afișează coordonatele x crescând spre dreapta, iar marginea dreaptă coordonatele y crescând în jos.} \\
\texttt{* Următoarele obiecte sunt plasate aleatoriu pe grila voastră: 'W', 'I', 'T', 'C', 'H'.} \\
\\ 
\\ 
\texttt{Celălalt jucător vede o variantă a tablei de joc, în care obiectele sunt plasate în diferite locații aleatorii. Nu puteți vedea tabla de joc a celuilalt jucător, iar el nu o poate vedea pe a dumneavoastră.} \\
\\ 
\\ 
\texttt{**Obiectiv:**} \\
\\ 
\texttt{Ambii jucători trebuie să mute obiectele de pe fundalul lor respectiv astfel încât obiectele identice să ajungă la aceleași coordonate. Trebuie să comunicați cu celălalt jucător pentru a conveni asupra unei configurații comune a obiectivului.} \\
\\ 
\\ 
\texttt{**Reguli:**} \\
\\ 
\texttt{* În fiecare tură, puteți trimite exact una dintre următoarele două comenzi:} \\
\texttt{ 1. `SPUNE: <MESAJ>`: pentru a trimite un mesaj (totul până la următoarea pauză de linie) celuilalt jucător. Eu îl voi transmite partenerului tău.} \\
\texttt{2. `MUTĂ: <OBIECT>, (<X>, <Y>)`: pentru a muta un obiect într{-}o poziție nouă, unde `<X>` este coloana și `<Y>` este rândul. Te voi informa dacă mutarea ta a avut succes sau nu.} \\
\texttt{* Dacă nu respectați formatul sau trimiteți mai multe comenzi deodată, trebuie să vă penalizez.} \\
\texttt{* Dacă ambii jucători acumulează mai mult de 12 penalități, amândoi pierdeți jocul.} \\
\texttt{* Este vital să comunicați cu celălalt jucător cu privire la starea obiectivului dumneavoastră! *Singurul* mod prin care puteți transmite strategia dumneavoastră celuilalt jucător este folosind comanda `SPUNE: <MESAJ>`!} \\
\\ 
\\ 
\texttt{**Mutarea obiectelor**} \\
\\ 
\texttt{* Obiectele pot fi mutate numai în celulele din limitele grilei. Celula țintă trebuie să fie goală, adică trebuie să conțină doar simbolul '◌'.} \\
\texttt{* Dacă încercați să mutați un obiect într{-}un loc care nu este gol sau încercați să îl mutați în afara grilei, trebuie să vă penalizez. Aveți dreptul la o nouă încercare.} \\
\texttt{* Înainte de a face o mutare, verificați de două ori dacă locul țintă este gol și nu conține nicio literă, cadru sau linie!} \\
\\ 
\\ 
\texttt{**Sfârșitul jocului**} \\
\\ 
\texttt{Dacă credeți că ați atins obiectivul de aliniere a tuturor obiectelor, puteți cere celuilalt jucător să încheie jocul trimițând `SPUNE: terminat?`. Dacă celălalt jucător vă cere să încheiați jocul, iar dvs. răspundeți cu `SPUNE: terminat!`, jocul se va încheia.} \\
\\ 
\texttt{Ambii jucători câștigă dacă jocul se încheie în decursul a 20 runde, unde o rundă este definită ca fiind trimiterea de către doi jucători a câte unei comenzi valide.} \\
\\ 
\texttt{**Punctaj:**} \\
\\ 
\texttt{Cu cât obiectele identice sunt mai apropiate pe ambele table de joc, cu atât primești mai multe puncte. Penalizările îți reduc punctele. Poți bate recordul?} \\
\\ 
\texttt{Celălalt jucător a început jocul trimițând acest mesaj:} \\
\\ 
\texttt{"Salut! Pe tabla mea, W este la (2,7), I la (6,5), T la (5,7), C la (1,6), H la (4,4). Unde sunt obiectele tale?"} \\
\\ 
\texttt{Care este prima ta comandă?} \\
            }
        }
    }
    & \\ \\

    \theutterance \stepcounter{utterance}  
    & & & \multicolumn{4}{p{0.6\linewidth}}{
        \cellcolor[rgb]{0.9,0.9,0.9}{
            \makecell[{{p{\linewidth}}}]{
                \texttt{\tiny{[GM$\langle$P2]}}
                \texttt{Salut! Să începem jocul colaborativ.} \\
\\ 
\texttt{Mai întâi, hai să identific pozițiile obiectelor pe tabla mea de joc:} \\
\\ 
\texttt{Din grila pe care o văd:} \\
\\ 
\texttt{{-} **W** este la (5,7)} \\
\texttt{{-} **I** este la (4,6)} \\
\texttt{{-} **T** este la (7,4)} \\
\texttt{{-} **C** este la (7,3)} \\
\texttt{{-} **H** este la (7,1)} \\
\\ 
\texttt{Acum, compar cu pozițiile de pe tabla celuilalt jucător:} \\
\\ 
\texttt{| Obiect | Tabla mea | Tabla partenerului |} \\
\texttt{|{-}{-}{-}{-}{-}{-}{-}{-}|{-}{-}{-}{-}{-}{-}{-}{-}{-}{-}{-}|{-}{-}{-}{-}{-}{-}{-}{-}{-}{-}{-}{-}{-}{-}{-}{-}{-}{-}{-}|} \\
\texttt{| W      | (5,7)     | (2,7)             |} \\
\texttt{| I      | (4,6)     | (6,5)             |} \\
\texttt{| T      | (7,4)     | (5,7)             |} \\
\texttt{| C      | (7,3)     | (1,6)             |} \\
\texttt{| H      | (7,1)     | (4,4)             |} \\
\\ 
\texttt{Îi comunic pozițiile obiectelor mele:} \\
\\ 
\texttt{```} \\
\texttt{SPUNE: Salut! Pe tabla mea: W=(5,7), I=(4,6), T=(7,4), C=(7,3), H=(7,1). Haide să cădem de acord asupra pozițiilor țintă. Ce zici de punctul central al grilei?} \\
\texttt{```} \\
            }
        }
    }
     \\ \\

    \theutterance \stepcounter{utterance}  
    & & & \multicolumn{2}{p{0.3\linewidth}}{
        \cellcolor[rgb]{0.9,0.9,0.9}{
            \makecell[{{p{\linewidth}}}]{
                \texttt{\tiny{[GM$|$GM]}}
                \texttt{Salut! Să începem jocul colaborativ.} \\
\\ 
\texttt{Mai întâi, hai să identific pozițiile obiectelor pe tabla mea de joc:} \\
\\ 
\texttt{Din grila pe care o văd:} \\
\\ 
\texttt{{-} **W** este la (5,7)} \\
\texttt{{-} **I** este la (4,6)} \\
\texttt{{-} **T** este la (7,4)} \\
\texttt{{-} **C** este la (7,3)} \\
\texttt{{-} **H** este la (7,1)} \\
\\ 
\texttt{Acum, compar cu pozițiile de pe tabla celuilalt jucător:} \\
\\ 
\texttt{| Obiect | Tabla mea | Tabla partenerului |} \\
\texttt{|{-}{-}{-}{-}{-}{-}{-}{-}|{-}{-}{-}{-}{-}{-}{-}{-}{-}{-}{-}|{-}{-}{-}{-}{-}{-}{-}{-}{-}{-}{-}{-}{-}{-}{-}{-}{-}{-}{-}|} \\
\texttt{| W      | (5,7)     | (2,7)             |} \\
\texttt{| I      | (4,6)     | (6,5)             |} \\
\texttt{| T      | (7,4)     | (5,7)             |} \\
\texttt{| C      | (7,3)     | (1,6)             |} \\
\texttt{| H      | (7,1)     | (4,4)             |} \\
\\ 
\texttt{Îi comunic pozițiile obiectelor mele:} \\
\\ 
\texttt{```} \\
\texttt{SPUNE: Salut! Pe tabla mea: W=(5,7), I=(4,6), T=(7,4), C=(7,3), H=(7,1). Haide să cădem de acord asupra pozițiilor țintă. Ce zici de punctul central al grilei?} \\
\texttt{```} \\
            }
        }
    }
    & & \\ \\

    \theutterance \stepcounter{utterance}  
    & & & \multicolumn{2}{p{0.3\linewidth}}{
        \cellcolor[rgb]{0.9,0.9,0.9}{
            \makecell[{{p{\linewidth}}}]{
                \texttt{\tiny{[GM$|$GM]}}
                \texttt{Invalid format: head is not empty: 'Salut! Să începem jocul colaborativ.} \\
\\ 
\texttt{Mai întâi, hai să identific pozițiile obiectelor pe tabla mea de joc:} \\
\\ 
\texttt{Din grila pe care o văd:} \\
\\ 
\texttt{{-} **W** este la (5,7)} \\
\texttt{{-} **I** este la (4,6)} \\
\texttt{{-} **T** este la (7,4)} \\
\texttt{{-} **C** este la (7,3)} \\
\texttt{{-} **H** este la (7,1)} \\
\\ 
\texttt{Acum, compar cu pozițiile de pe tabla celuilalt jucător:} \\
\\ 
\texttt{| Obiect | Tabla mea | Tabla partenerului |} \\
\texttt{|{-}{-}{-}{-}{-}{-}{-}{-}|{-}{-}{-}{-}{-}{-}{-}{-}{-}{-}{-}|{-}{-}{-}{-}{-}{-}{-}{-}{-}{-}{-}{-}{-}{-}{-}{-}{-}{-}{-}|} \\
\texttt{| W      | (5,7)     | (2,7)             |} \\
\texttt{| I      | (4,6)     | (6,5)             |} \\
\texttt{| T      | (7,4)     | (5,7)             |} \\
\texttt{| C      | (7,3)     | (1,6)             |} \\
\texttt{| H      | (7,1)     | (4,4)             |} \\
\\ 
\texttt{Îi comunic pozițiile obiectelor mele:} \\
\\ 
\\ 
\texttt{'} \\
            }
        }
    }
    & & \\ \\

    \theutterance \stepcounter{utterance}  
    & & \multicolumn{4}{p{0.6\linewidth}}{
        \cellcolor[rgb]{0.9,0.9,0.9}{
            \makecell[{{p{\linewidth}}}]{
                \texttt{\tiny{[GM$\rangle$P2]}}
                \texttt{Penalizare: Mesajul dvs. nu trebuie să conțină nimic înaintea comenzii!} \\
\\ 
\texttt{Asigurați{-}vă că răspunsul dvs. conține doar `SPUNE: <MESAJ>` sau `MUTĂ: <OBIECT>, (<X>, <Y>)`, și nimic altceva!} \\
\\ 
\texttt{Ați acumulat în mod colectiv 2 din 12 penalizări.} \\
\texttt{Vă rugăm să încercați din nou!} \\
            }
        }
    }
    & \\ \\

    \theutterance \stepcounter{utterance}  
    & & & \multicolumn{4}{p{0.6\linewidth}}{
        \cellcolor[rgb]{0.9,0.9,0.9}{
            \makecell[{{p{\linewidth}}}]{
                \texttt{\tiny{[GM$\langle$P2]}}
                \texttt{SPUNE: Pe tabla mea: W=(5,7), I=(4,6), T=(7,4), C=(7,3), H=(7,1). Hai să alegem poziții comune. Propun să ne întâlnim la centru.} \\
            }
        }
    }
     \\ \\

    \theutterance \stepcounter{utterance}  
    & & & \multicolumn{2}{p{0.3\linewidth}}{
        \cellcolor[rgb]{0.9,0.9,0.9}{
            \makecell[{{p{\linewidth}}}]{
                \texttt{\tiny{[GM$|$GM]}}
                \texttt{SPUNE: Pe tabla mea: W=(5,7), I=(4,6), T=(7,4), C=(7,3), H=(7,1). Hai să alegem poziții comune. Propun să ne întâlnim la centru.} \\
            }
        }
    }
    & & \\ \\

    \theutterance \stepcounter{utterance}  
    & & \multicolumn{4}{p{0.6\linewidth}}{
        \cellcolor[rgb]{0.9,0.9,0.9}{
            \makecell[{{p{\linewidth}}}]{
                \texttt{\tiny{[P1$\langle$GM]}}
                \texttt{Mesajul dvs. a fost transmis celuilalt jucător.} \\
\\ 
\\ 
\texttt{În prezent jucați runda 1 din maximum 20.} \\
\texttt{Ați acumulat în mod colectiv 2 din 12 penalizări.} \\
\texttt{Celălalt jucător a trimis acest mesaj:} \\
\\ 
\\ 
\texttt{"Pe tabla mea: W=(5,7), I=(4,6), T=(7,4), C=(7,3), H=(7,1). Hai să alegem poziții comune. Propun să ne întâlnim la centru."} \\
\\ 
\\ 
\texttt{Care este următoarea dvs. comandă?} \\
            }
        }
    }
    & \\ \\

    \theutterance \stepcounter{utterance}  
    & \multicolumn{4}{p{0.6\linewidth}}{
        \cellcolor[rgb]{0.9,0.9,0.9}{
            \makecell[{{p{\linewidth}}}]{
                \texttt{\tiny{[P1$\rangle$GM]}}
                \texttt{SPUNE: Am înțeles pozițiile tale. Propun să mutăm W la (4,4), I la (4,4) nu merge pentru că e ocupat. Hai să alegem: W la (3,3), I la (3,4), T la (3,5), C la (3,6), H la (3,7). Ce zici?} \\
            }
        }
    }
    & & \\ \\

    \theutterance \stepcounter{utterance}  
    & & & \multicolumn{2}{p{0.3\linewidth}}{
        \cellcolor[rgb]{0.9,0.9,0.9}{
            \makecell[{{p{\linewidth}}}]{
                \texttt{\tiny{[GM$|$GM]}}
                \texttt{SPUNE: Am înțeles pozițiile tale. Propun să mutăm W la (4,4), I la (4,4) nu merge pentru că e ocupat. Hai să alegem: W la (3,3), I la (3,4), T la (3,5), C la (3,6), H la (3,7). Ce zici?} \\
            }
        }
    }
    & & \\ \\

    \theutterance \stepcounter{utterance}  
    & & \multicolumn{4}{p{0.6\linewidth}}{
        \cellcolor[rgb]{0.9,0.9,0.9}{
            \makecell[{{p{\linewidth}}}]{
                \texttt{\tiny{[GM$\rangle$P2]}}
                \texttt{Mesajul dvs. a fost transmis celuilalt jucător.} \\
\\ 
\\ 
\texttt{În prezent jucați runda 2 din maximum 20.} \\
\texttt{Ați acumulat în mod colectiv 2 din 12 penalizări.} \\
\texttt{Celălalt jucător a trimis acest mesaj:} \\
\\ 
\\ 
\texttt{"Am înțeles pozițiile tale. Propun să mutăm W la (4,4), I la (4,4) nu merge pentru că e ocupat. Hai să alegem: W la (3,3), I la (3,4), T la (3,5), C la (3,6), H la (3,7). Ce zici?"} \\
\\ 
\\ 
\texttt{Care este următoarea dvs. comandă?} \\
            }
        }
    }
    & \\ \\

    \theutterance \stepcounter{utterance}  
    & & & \multicolumn{4}{p{0.6\linewidth}}{
        \cellcolor[rgb]{0.9,0.9,0.9}{
            \makecell[{{p{\linewidth}}}]{
                \texttt{\tiny{[GM$\langle$P2]}}
                \texttt{SPUNE: Interesantă propunerea! Dar verificați coloana 3 {-} în grila mea sunt multe obstacole acolo (linii de cadru). Ce ziceți de o coloană mai accesibilă, precum coloana 4? Propun: W=(4,3), I=(4,4), T=(4,5), C=(4,6), H=(4,7). Verificați dacă aceste poziții sunt libere la tine.} \\
            }
        }
    }
     \\ \\

    \theutterance \stepcounter{utterance}  
    & & & \multicolumn{2}{p{0.3\linewidth}}{
        \cellcolor[rgb]{0.9,0.9,0.9}{
            \makecell[{{p{\linewidth}}}]{
                \texttt{\tiny{[GM$|$GM]}}
                \texttt{SPUNE: Interesantă propunerea! Dar verificați coloana 3 {-} în grila mea sunt multe obstacole acolo (linii de cadru). Ce ziceți de o coloană mai accesibilă, precum coloana 4? Propun: W=(4,3), I=(4,4), T=(4,5), C=(4,6), H=(4,7). Verificați dacă aceste poziții sunt libere la tine.} \\
            }
        }
    }
    & & \\ \\

    \theutterance \stepcounter{utterance}  
    & & \multicolumn{4}{p{0.6\linewidth}}{
        \cellcolor[rgb]{0.9,0.9,0.9}{
            \makecell[{{p{\linewidth}}}]{
                \texttt{\tiny{[P1$\langle$GM]}}
                \texttt{Mesajul dvs. a fost transmis celuilalt jucător.} \\
\\ 
\\ 
\texttt{În prezent jucați runda 2 din maximum 20.} \\
\texttt{Ați acumulat în mod colectiv 2 din 12 penalizări.} \\
\texttt{Celălalt jucător a trimis acest mesaj:} \\
\\ 
\\ 
\texttt{"Interesantă propunerea! Dar verificați coloana 3 {-} în grila mea sunt multe obstacole acolo (linii de cadru). Ce ziceți de o coloană mai accesibilă, precum coloana 4? Propun: W=(4,3), I=(4,4), T=(4,5), C=(4,6), H=(4,7). Verificați dacă aceste poziții sunt libere la tine."} \\
\\ 
\\ 
\texttt{Care este următoarea dvs. comandă?} \\
            }
        }
    }
    & \\ \\

    \theutterance \stepcounter{utterance}  
    & \multicolumn{4}{p{0.6\linewidth}}{
        \cellcolor[rgb]{0.9,0.9,0.9}{
            \makecell[{{p{\linewidth}}}]{
                \texttt{\tiny{[P1$\rangle$GM]}}
                \texttt{SPUNE: Pe tabla mea, coloana 4 are următoarele: (4,1) e liber, (4,2) e liber, (4,3) e liber, (4,4) e ocupat de H, (4,5) e ocupat de I?, stai să verific. De fapt, H e la (4,4), deci (4,4) nu e liber. Propun coloana 5: W=(5,2), I=(5,3), T=(5,4), C=(5,5), H=(5,6). Sunt libere la tine?} \\
            }
        }
    }
    & & \\ \\

    \theutterance \stepcounter{utterance}  
    & & & \multicolumn{2}{p{0.3\linewidth}}{
        \cellcolor[rgb]{0.9,0.9,0.9}{
            \makecell[{{p{\linewidth}}}]{
                \texttt{\tiny{[GM$|$GM]}}
                \texttt{SPUNE: Pe tabla mea, coloana 4 are următoarele: (4,1) e liber, (4,2) e liber, (4,3) e liber, (4,4) e ocupat de H, (4,5) e ocupat de I?, stai să verific. De fapt, H e la (4,4), deci (4,4) nu e liber. Propun coloana 5: W=(5,2), I=(5,3), T=(5,4), C=(5,5), H=(5,6). Sunt libere la tine?} \\
            }
        }
    }
    & & \\ \\

    \theutterance \stepcounter{utterance}  
    & & \multicolumn{4}{p{0.6\linewidth}}{
        \cellcolor[rgb]{0.9,0.9,0.9}{
            \makecell[{{p{\linewidth}}}]{
                \texttt{\tiny{[GM$\rangle$P2]}}
                \texttt{Mesajul dvs. a fost transmis celuilalt jucător.} \\
\\ 
\\ 
\texttt{În prezent jucați runda 3 din maximum 20.} \\
\texttt{Ați acumulat în mod colectiv 2 din 12 penalizări.} \\
\texttt{Celălalt jucător a trimis acest mesaj:} \\
\\ 
\\ 
\texttt{"Pe tabla mea, coloana 4 are următoarele: (4,1) e liber, (4,2) e liber, (4,3) e liber, (4,4) e ocupat de H, (4,5) e ocupat de I?, stai să verific. De fapt, H e la (4,4), deci (4,4) nu e liber. Propun coloana 5: W=(5,2), I=(5,3), T=(5,4), C=(5,5), H=(5,6). Sunt libere la tine?"} \\
\\ 
\\ 
\texttt{Care este următoarea dvs. comandă?} \\
            }
        }
    }
    & \\ \\

    \theutterance \stepcounter{utterance}  
    & & & \multicolumn{4}{p{0.6\linewidth}}{
        \cellcolor[rgb]{0.9,0.9,0.9}{
            \makecell[{{p{\linewidth}}}]{
                \texttt{\tiny{[GM$\langle$P2]}}
                \texttt{SPUNE: Verific acum coloana 5 pe tabla mea: (5,1) e liber, (5,2) e liber, (5,3) e liber, (5,4) e liber, (5,5) e liber, (5,6) e liber, (5,7) e ocupat de W momentan. Dar W se va muta, deci da, coloana 5 funcționează! Accept propunerea: W=(5,2), I=(5,3), T=(5,4), C=(5,5), H=(5,6). Începem mutările?} \\
            }
        }
    }
     \\ \\

    \theutterance \stepcounter{utterance}  
    & & & \multicolumn{2}{p{0.3\linewidth}}{
        \cellcolor[rgb]{0.9,0.9,0.9}{
            \makecell[{{p{\linewidth}}}]{
                \texttt{\tiny{[GM$|$GM]}}
                \texttt{SPUNE: Verific acum coloana 5 pe tabla mea: (5,1) e liber, (5,2) e liber, (5,3) e liber, (5,4) e liber, (5,5) e liber, (5,6) e liber, (5,7) e ocupat de W momentan. Dar W se va muta, deci da, coloana 5 funcționează! Accept propunerea: W=(5,2), I=(5,3), T=(5,4), C=(5,5), H=(5,6). Începem mutările?} \\
            }
        }
    }
    & & \\ \\

    \theutterance \stepcounter{utterance}  
    & & \multicolumn{4}{p{0.6\linewidth}}{
        \cellcolor[rgb]{0.9,0.9,0.9}{
            \makecell[{{p{\linewidth}}}]{
                \texttt{\tiny{[P1$\langle$GM]}}
                \texttt{Mesajul dvs. a fost transmis celuilalt jucător.} \\
\\ 
\\ 
\texttt{În prezent jucați runda 3 din maximum 20.} \\
\texttt{Ați acumulat în mod colectiv 2 din 12 penalizări.} \\
\texttt{Celălalt jucător a trimis acest mesaj:} \\
\\ 
\\ 
\texttt{"Verific acum coloana 5 pe tabla mea: (5,1) e liber, (5,2) e liber, (5,3) e liber, (5,4) e liber, (5,5) e liber, (5,6) e liber, (5,7) e ocupat de W momentan. Dar W se va muta, deci da, coloana 5 funcționează! Accept propunerea: W=(5,2), I=(5,3), T=(5,4), C=(5,5), H=(5,6). Începem mutările?"} \\
\\ 
\\ 
\texttt{Care este următoarea dvs. comandă?} \\
            }
        }
    }
    & \\ \\

    \theutterance \stepcounter{utterance}  
    & \multicolumn{4}{p{0.6\linewidth}}{
        \cellcolor[rgb]{0.9,0.9,0.9}{
            \makecell[{{p{\linewidth}}}]{
                \texttt{\tiny{[P1$\rangle$GM]}}
                \texttt{SPUNE: Da, începem mutările! Încep cu W. Mut W la (5,2).} \\
            }
        }
    }
    & & \\ \\

    \theutterance \stepcounter{utterance}  
    & & & \multicolumn{2}{p{0.3\linewidth}}{
        \cellcolor[rgb]{0.9,0.9,0.9}{
            \makecell[{{p{\linewidth}}}]{
                \texttt{\tiny{[GM$|$GM]}}
                \texttt{SPUNE: Da, începem mutările! Încep cu W. Mut W la (5,2).} \\
            }
        }
    }
    & & \\ \\

    \theutterance \stepcounter{utterance}  
    & & \multicolumn{4}{p{0.6\linewidth}}{
        \cellcolor[rgb]{0.9,0.9,0.9}{
            \makecell[{{p{\linewidth}}}]{
                \texttt{\tiny{[GM$\rangle$P2]}}
                \texttt{Mesajul dvs. a fost transmis celuilalt jucător.} \\
\\ 
\\ 
\texttt{În prezent jucați runda 4 din maximum 20.} \\
\texttt{Ați acumulat în mod colectiv 2 din 12 penalizări.} \\
\texttt{Celălalt jucător a trimis acest mesaj:} \\
\\ 
\\ 
\texttt{"Da, începem mutările! Încep cu W. Mut W la (5,2)."} \\
\\ 
\\ 
\texttt{Care este următoarea dvs. comandă?} \\
            }
        }
    }
    & \\ \\

    \theutterance \stepcounter{utterance}  
    & & & \multicolumn{4}{p{0.6\linewidth}}{
        \cellcolor[rgb]{0.9,0.9,0.9}{
            \makecell[{{p{\linewidth}}}]{
                \texttt{\tiny{[GM$\langle$P2]}}
                \texttt{MUTĂ: W, (5,2)} \\
            }
        }
    }
     \\ \\

    \theutterance \stepcounter{utterance}  
    & & & \multicolumn{2}{p{0.3\linewidth}}{
        \cellcolor[rgb]{0.9,0.9,0.9}{
            \makecell[{{p{\linewidth}}}]{
                \texttt{\tiny{[GM$|$GM]}}
                \texttt{MUTĂ: W, (5,2)} \\
            }
        }
    }
    & & \\ \\

    \theutterance \stepcounter{utterance}  
    & & & \multicolumn{2}{p{0.3\linewidth}}{
        \cellcolor[rgb]{0.9,0.9,0.9}{
            \makecell[{{p{\linewidth}}}]{
                \texttt{\tiny{[GM$|$GM]}}
                \texttt{Penalizare: (5,2) nu este gol, ci conține '┬'.} \\
\\ 
\texttt{Ați acumulat în mod colectiv 3 din 12 penalizări.} \\
\texttt{Vă rugăm să încercați din nou!} \\
            }
        }
    }
    & & \\ \\

    \theutterance \stepcounter{utterance}  
    & & \multicolumn{4}{p{0.6\linewidth}}{
        \cellcolor[rgb]{0.9,0.9,0.9}{
            \makecell[{{p{\linewidth}}}]{
                \texttt{\tiny{[GM$\rangle$P2]}}
                \texttt{Penalizare: (5,2) nu este gol, ci conține '┬'.} \\
\\ 
\texttt{Ați acumulat în mod colectiv 3 din 12 penalizări.} \\
\texttt{Vă rugăm să încercați din nou!} \\
            }
        }
    }
    & \\ \\

    \theutterance \stepcounter{utterance}  
    & & & \multicolumn{4}{p{0.6\linewidth}}{
        \cellcolor[rgb]{0.9,0.9,0.9}{
            \makecell[{{p{\linewidth}}}]{
                \texttt{\tiny{[GM$\langle$P2]}}
                \texttt{Îmi pare rău pentru eroare. Să verific cu atenție grila.} \\
\\ 
\texttt{Poziția (5,2) în grila mea {-} mă uit la rândul 2, coloana 5. Rândul 2 arată: `╟─┼──┬┘◌║`, deci coloana 5 conține `┬`, un caracter de cadru, nu `◌`.} \\
\\ 
\texttt{Atunci coloana 5 rândul 2 nu e liberă pe tabla mea! Trebuie să verific toate pozițiile propuse:} \\
\\ 
\texttt{{-} (5,2): rândul 2, `╟─┼──┬┘◌║` → coloana 5 = `┬` (cadru)} \\
\texttt{{-} (5,3): rândul 3, `║◌│◌◌│◌C║` → coloana 5 = `│` (cadru)} \\
\texttt{{-} (5,4): rândul 4, `╟─┴──┘◌T║` → coloana 5 = `┘` (cadru)} \\
\texttt{{-} (5,5): rândul 5, `║◌◌◌◌◌◌◌║` → coloana 5 = `◌` (liber)} \\
\texttt{{-} (5,6): rândul 6, `║◌◌◌I◌◌◌║` → coloana 5 = `◌` (liber)} \\
\\ 
\texttt{Deci coloana 5 are probleme! Trebuie să anunț partenerul.} \\
\\ 
\texttt{SPUNE: Stop! Am verificat mai atent grila mea. Coloana 5 la mine: (5,2) are un cadru '┬', (5,3) are '│', (5,4) are '┘'. Doar (5,5) și (5,6) sunt libere. Coloana 5 nu merge! Hai să găsim altă coloană. Verific coloana 6 la mine: (6,1) liber, (6,2) liber, (6,3) liber, (6,4) liber, (6,5) ocupat de I, (6,6) liber, (6,7) liber. Coloana 6 e mai bună. Propun: W=(6,1), I=(6,2), T=(6,3), C=(6,4), H=(6,6). Verifică la tine!} \\
            }
        }
    }
     \\ \\

    \theutterance \stepcounter{utterance}  
    & & & \multicolumn{2}{p{0.3\linewidth}}{
        \cellcolor[rgb]{0.9,0.9,0.9}{
            \makecell[{{p{\linewidth}}}]{
                \texttt{\tiny{[GM$|$GM]}}
                \texttt{Îmi pare rău pentru eroare. Să verific cu atenție grila.} \\
\\ 
\texttt{Poziția (5,2) în grila mea {-} mă uit la rândul 2, coloana 5. Rândul 2 arată: `╟─┼──┬┘◌║`, deci coloana 5 conține `┬`, un caracter de cadru, nu `◌`.} \\
\\ 
\texttt{Atunci coloana 5 rândul 2 nu e liberă pe tabla mea! Trebuie să verific toate pozițiile propuse:} \\
\\ 
\texttt{{-} (5,2): rândul 2, `╟─┼──┬┘◌║` → coloana 5 = `┬` (cadru)} \\
\texttt{{-} (5,3): rândul 3, `║◌│◌◌│◌C║` → coloana 5 = `│` (cadru)} \\
\texttt{{-} (5,4): rândul 4, `╟─┴──┘◌T║` → coloana 5 = `┘` (cadru)} \\
\texttt{{-} (5,5): rândul 5, `║◌◌◌◌◌◌◌║` → coloana 5 = `◌` (liber)} \\
\texttt{{-} (5,6): rândul 6, `║◌◌◌I◌◌◌║` → coloana 5 = `◌` (liber)} \\
\\ 
\texttt{Deci coloana 5 are probleme! Trebuie să anunț partenerul.} \\
\\ 
\texttt{SPUNE: Stop! Am verificat mai atent grila mea. Coloana 5 la mine: (5,2) are un cadru '┬', (5,3) are '│', (5,4) are '┘'. Doar (5,5) și (5,6) sunt libere. Coloana 5 nu merge! Hai să găsim altă coloană. Verific coloana 6 la mine: (6,1) liber, (6,2) liber, (6,3) liber, (6,4) liber, (6,5) ocupat de I, (6,6) liber, (6,7) liber. Coloana 6 e mai bună. Propun: W=(6,1), I=(6,2), T=(6,3), C=(6,4), H=(6,6). Verifică la tine!} \\
            }
        }
    }
    & & \\ \\

    \theutterance \stepcounter{utterance}  
    & & & \multicolumn{2}{p{0.3\linewidth}}{
        \cellcolor[rgb]{0.9,0.9,0.9}{
            \makecell[{{p{\linewidth}}}]{
                \texttt{\tiny{[GM$|$GM]}}
                \texttt{Invalid format: head is not empty: 'Îmi pare rău pentru eroare. Să verific cu atenție grila.} \\
\\ 
\texttt{Poziția (5,2) în grila mea {-} mă uit la rândul 2, coloana 5. Rândul 2 arată: ╟─┼──┬┘◌║, deci coloana 5 conține ┬, un caracter de cadru, nu ◌.} \\
\\ 
\texttt{Atunci coloana 5 rândul 2 nu e liberă pe tabla mea! Trebuie să verific toate pozițiile propuse:} \\
\\ 
\texttt{{-} (5,2): rândul 2, ╟─┼──┬┘◌║ → coloana 5 = ┬ (cadru)} \\
\texttt{{-} (5,3): rândul 3, ║◌│◌◌│◌C║ → coloana 5 = │ (cadru)} \\
\texttt{{-} (5,4): rândul 4, ╟─┴──┘◌T║ → coloana 5 = ┘ (cadru)} \\
\texttt{{-} (5,5): rândul 5, ║◌◌◌◌◌◌◌║ → coloana 5 = ◌ (liber)} \\
\texttt{{-} (5,6): rândul 6, ║◌◌◌I◌◌◌║ → coloana 5 = ◌ (liber)} \\
\\ 
\texttt{Deci coloana 5 are probleme! Trebuie să anunț partenerul.} \\
\\ 
\texttt{'} \\
            }
        }
    }
    & & \\ \\

    \theutterance \stepcounter{utterance}  
    & & \multicolumn{4}{p{0.6\linewidth}}{
        \cellcolor[rgb]{0.9,0.9,0.9}{
            \makecell[{{p{\linewidth}}}]{
                \texttt{\tiny{[GM$\rangle$P2]}}
                \texttt{Penalizare: Mesajul dvs. nu trebuie să conțină nimic înaintea comenzii!} \\
\\ 
\texttt{Asigurați{-}vă că răspunsul dvs. conține doar `SPUNE: <MESAJ>` sau `MUTĂ: <OBIECT>, (<X>, <Y>)`, și nimic altceva!} \\
\\ 
\texttt{Ați acumulat în mod colectiv 4 din 12 penalizări.} \\
\texttt{Vă rugăm să încercați din nou!} \\
            }
        }
    }
    & \\ \\

    \theutterance \stepcounter{utterance}  
    & & & \multicolumn{4}{p{0.6\linewidth}}{
        \cellcolor[rgb]{0.9,0.9,0.9}{
            \makecell[{{p{\linewidth}}}]{
                \texttt{\tiny{[GM$\langle$P2]}}
                \texttt{SPUNE: Stop! Am verificat mai atent grila mea. Coloana 5 la mine: (5,2) are un cadru, (5,3) are cadru, (5,4) are cadru. Doar (5,5) și (5,6) sunt libere. Coloana 5 nu merge! Propun coloana 6: W=(6,1), I=(6,2), T=(6,3), C=(6,4), H=(6,6). Verifică la tine!} \\
            }
        }
    }
     \\ \\

    \theutterance \stepcounter{utterance}  
    & & & \multicolumn{2}{p{0.3\linewidth}}{
        \cellcolor[rgb]{0.9,0.9,0.9}{
            \makecell[{{p{\linewidth}}}]{
                \texttt{\tiny{[GM$|$GM]}}
                \texttt{SPUNE: Stop! Am verificat mai atent grila mea. Coloana 5 la mine: (5,2) are un cadru, (5,3) are cadru, (5,4) are cadru. Doar (5,5) și (5,6) sunt libere. Coloana 5 nu merge! Propun coloana 6: W=(6,1), I=(6,2), T=(6,3), C=(6,4), H=(6,6). Verifică la tine!} \\
            }
        }
    }
    & & \\ \\

    \theutterance \stepcounter{utterance}  
    & & \multicolumn{4}{p{0.6\linewidth}}{
        \cellcolor[rgb]{0.9,0.9,0.9}{
            \makecell[{{p{\linewidth}}}]{
                \texttt{\tiny{[P1$\langle$GM]}}
                \texttt{Mesajul dvs. a fost transmis celuilalt jucător.} \\
\\ 
\\ 
\texttt{În prezent jucați runda 4 din maximum 20.} \\
\texttt{Ați acumulat în mod colectiv 4 din 12 penalizări.} \\
\texttt{Celălalt jucător a trimis acest mesaj:} \\
\\ 
\\ 
\texttt{"Stop! Am verificat mai atent grila mea. Coloana 5 la mine: (5,2) are un cadru, (5,3) are cadru, (5,4) are cadru. Doar (5,5) și (5,6) sunt libere. Coloana 5 nu merge! Propun coloana 6: W=(6,1), I=(6,2), T=(6,3), C=(6,4), H=(6,6). Verifică la tine!"} \\
\\ 
\\ 
\texttt{Care este următoarea dvs. comandă?} \\
            }
        }
    }
    & \\ \\

    \theutterance \stepcounter{utterance}  
    & \multicolumn{4}{p{0.6\linewidth}}{
        \cellcolor[rgb]{0.9,0.9,0.9}{
            \makecell[{{p{\linewidth}}}]{
                \texttt{\tiny{[P1$\rangle$GM]}}
                \texttt{SPUNE: Verific coloana 6 pe tabla mea: (6,1) e liber, (6,2) e liber, (6,3) e liber, (6,4) e liber, (6,5) e ocupat de I, (6,6) e liber, (6,7) e liber. Dacă mut I de la (6,5), atunci (6,6) e liber. Dar tu ai zis H=(6,6), nu e problemă. Accept coloana 6: W=(6,1), I=(6,2), T=(6,3), C=(6,4), H=(6,6). Confirmat?} \\
            }
        }
    }
    & & \\ \\

    \theutterance \stepcounter{utterance}  
    & & & \multicolumn{2}{p{0.3\linewidth}}{
        \cellcolor[rgb]{0.9,0.9,0.9}{
            \makecell[{{p{\linewidth}}}]{
                \texttt{\tiny{[GM$|$GM]}}
                \texttt{SPUNE: Verific coloana 6 pe tabla mea: (6,1) e liber, (6,2) e liber, (6,3) e liber, (6,4) e liber, (6,5) e ocupat de I, (6,6) e liber, (6,7) e liber. Dacă mut I de la (6,5), atunci (6,6) e liber. Dar tu ai zis H=(6,6), nu e problemă. Accept coloana 6: W=(6,1), I=(6,2), T=(6,3), C=(6,4), H=(6,6). Confirmat?} \\
            }
        }
    }
    & & \\ \\

    \theutterance \stepcounter{utterance}  
    & & \multicolumn{4}{p{0.6\linewidth}}{
        \cellcolor[rgb]{0.9,0.9,0.9}{
            \makecell[{{p{\linewidth}}}]{
                \texttt{\tiny{[GM$\rangle$P2]}}
                \texttt{Mesajul dvs. a fost transmis celuilalt jucător.} \\
\\ 
\\ 
\texttt{În prezent jucați runda 5 din maximum 20.} \\
\texttt{Ați acumulat în mod colectiv 4 din 12 penalizări.} \\
\texttt{Celălalt jucător a trimis acest mesaj:} \\
\\ 
\\ 
\texttt{"Verific coloana 6 pe tabla mea: (6,1) e liber, (6,2) e liber, (6,3) e liber, (6,4) e liber, (6,5) e ocupat de I, (6,6) e liber, (6,7) e liber. Dacă mut I de la (6,5), atunci (6,6) e liber. Dar tu ai zis H=(6,6), nu e problemă. Accept coloana 6: W=(6,1), I=(6,2), T=(6,3), C=(6,4), H=(6,6). Confirmat?"} \\
\\ 
\\ 
\texttt{Care este următoarea dvs. comandă?} \\
            }
        }
    }
    & \\ \\

    \theutterance \stepcounter{utterance}  
    & & & \multicolumn{4}{p{0.6\linewidth}}{
        \cellcolor[rgb]{0.9,0.9,0.9}{
            \makecell[{{p{\linewidth}}}]{
                \texttt{\tiny{[GM$\langle$P2]}}
                \texttt{SPUNE: Confirmat! Începem mutările. Mut primul W la (6,1).} \\
            }
        }
    }
     \\ \\

    \theutterance \stepcounter{utterance}  
    & & & \multicolumn{2}{p{0.3\linewidth}}{
        \cellcolor[rgb]{0.9,0.9,0.9}{
            \makecell[{{p{\linewidth}}}]{
                \texttt{\tiny{[GM$|$GM]}}
                \texttt{SPUNE: Confirmat! Începem mutările. Mut primul W la (6,1).} \\
            }
        }
    }
    & & \\ \\

    \theutterance \stepcounter{utterance}  
    & & \multicolumn{4}{p{0.6\linewidth}}{
        \cellcolor[rgb]{0.9,0.9,0.9}{
            \makecell[{{p{\linewidth}}}]{
                \texttt{\tiny{[P1$\langle$GM]}}
                \texttt{Mesajul dvs. a fost transmis celuilalt jucător.} \\
\\ 
\\ 
\texttt{În prezent jucați runda 5 din maximum 20.} \\
\texttt{Ați acumulat în mod colectiv 4 din 12 penalizări.} \\
\texttt{Celălalt jucător a trimis acest mesaj:} \\
\\ 
\\ 
\texttt{"Confirmat! Începem mutările. Mut primul W la (6,1)."} \\
\\ 
\\ 
\texttt{Care este următoarea dvs. comandă?} \\
            }
        }
    }
    & \\ \\

    \theutterance \stepcounter{utterance}  
    & \multicolumn{4}{p{0.6\linewidth}}{
        \cellcolor[rgb]{0.9,0.9,0.9}{
            \makecell[{{p{\linewidth}}}]{
                \texttt{\tiny{[P1$\rangle$GM]}}
                \texttt{MUTĂ: W, (6, 1)} \\
            }
        }
    }
    & & \\ \\

    \theutterance \stepcounter{utterance}  
    & & & \multicolumn{2}{p{0.3\linewidth}}{
        \cellcolor[rgb]{0.9,0.9,0.9}{
            \makecell[{{p{\linewidth}}}]{
                \texttt{\tiny{[GM$|$GM]}}
                \texttt{MUTĂ: W, (6, 1)} \\
            }
        }
    }
    & & \\ \\

    \theutterance \stepcounter{utterance}  
    & & & \multicolumn{2}{p{0.3\linewidth}}{
        \cellcolor[rgb]{0.9,0.9,0.9}{
            \makecell[{{p{\linewidth}}}]{
                \texttt{\tiny{[GM$|$GM]}}
                \texttt{Penalizare: (6,1) nu este gol, ci conține '│'.} \\
\\ 
\texttt{Ați acumulat în mod colectiv 5 din 12 penalizări.} \\
\texttt{Vă rugăm să încercați din nou!} \\
            }
        }
    }
    & & \\ \\

    \theutterance \stepcounter{utterance}  
    & & \multicolumn{4}{p{0.6\linewidth}}{
        \cellcolor[rgb]{0.9,0.9,0.9}{
            \makecell[{{p{\linewidth}}}]{
                \texttt{\tiny{[P1$\langle$GM]}}
                \texttt{Penalizare: (6,1) nu este gol, ci conține '│'.} \\
\\ 
\texttt{Ați acumulat în mod colectiv 5 din 12 penalizări.} \\
\texttt{Vă rugăm să încercați din nou!} \\
            }
        }
    }
    & \\ \\

    \theutterance \stepcounter{utterance}  
    & \multicolumn{4}{p{0.6\linewidth}}{
        \cellcolor[rgb]{0.9,0.9,0.9}{
            \makecell[{{p{\linewidth}}}]{
                \texttt{\tiny{[P1$\rangle$GM]}}
                \texttt{MUTĂ: W, (6, 2)} \\
            }
        }
    }
    & & \\ \\

    \theutterance \stepcounter{utterance}  
    & & & \multicolumn{2}{p{0.3\linewidth}}{
        \cellcolor[rgb]{0.9,0.9,0.9}{
            \makecell[{{p{\linewidth}}}]{
                \texttt{\tiny{[GM$|$GM]}}
                \texttt{MUTĂ: W, (6, 2)} \\
            }
        }
    }
    & & \\ \\

    \theutterance \stepcounter{utterance}  
    & & & \multicolumn{2}{p{0.3\linewidth}}{
        \cellcolor[rgb]{0.9,0.9,0.9}{
            \makecell[{{p{\linewidth}}}]{
                \texttt{\tiny{[GM$|$GM]}}
                \texttt{Penalizare: (6,2) nu este gol, ci conține '┘'.} \\
\\ 
\texttt{Ați acumulat în mod colectiv 6 din 12 penalizări.} \\
\texttt{Vă rugăm să încercați din nou!} \\
            }
        }
    }
    & & \\ \\

    \theutterance \stepcounter{utterance}  
    & & \multicolumn{4}{p{0.6\linewidth}}{
        \cellcolor[rgb]{0.9,0.9,0.9}{
            \makecell[{{p{\linewidth}}}]{
                \texttt{\tiny{[P1$\langle$GM]}}
                \texttt{Penalizare: (6,2) nu este gol, ci conține '┘'.} \\
\\ 
\texttt{Ați acumulat în mod colectiv 6 din 12 penalizări.} \\
\texttt{Vă rugăm să încercați din nou!} \\
            }
        }
    }
    & \\ \\

    \theutterance \stepcounter{utterance}  
    & \multicolumn{4}{p{0.6\linewidth}}{
        \cellcolor[rgb]{0.9,0.9,0.9}{
            \makecell[{{p{\linewidth}}}]{
                \texttt{\tiny{[P1$\rangle$GM]}}
                \texttt{SPUNE: Pe tabla mea, coloana 6 la rândurile 1{-}4 conține caractere de cadru: (6,1) e '│', (6,2) e '┘', (6,3) e '┼', (6,4) e '┌'. Nu pot muta acolo. Coloana 6 liberă e doar (6,6) și (6,7). Trebuie să găsim altă coloană. Cum arată coloana 7 la tine?} \\
            }
        }
    }
    & & \\ \\

    \theutterance \stepcounter{utterance}  
    & & & \multicolumn{2}{p{0.3\linewidth}}{
        \cellcolor[rgb]{0.9,0.9,0.9}{
            \makecell[{{p{\linewidth}}}]{
                \texttt{\tiny{[GM$|$GM]}}
                \texttt{SPUNE: Pe tabla mea, coloana 6 la rândurile 1{-}4 conține caractere de cadru: (6,1) e '│', (6,2) e '┘', (6,3) e '┼', (6,4) e '┌'. Nu pot muta acolo. Coloana 6 liberă e doar (6,6) și (6,7). Trebuie să găsim altă coloană. Cum arată coloana 7 la tine?} \\
            }
        }
    }
    & & \\ \\

    \theutterance \stepcounter{utterance}  
    & & \multicolumn{4}{p{0.6\linewidth}}{
        \cellcolor[rgb]{0.9,0.9,0.9}{
            \makecell[{{p{\linewidth}}}]{
                \texttt{\tiny{[GM$\rangle$P2]}}
                \texttt{Mesajul dvs. a fost transmis celuilalt jucător.} \\
\\ 
\\ 
\texttt{În prezent jucați runda 6 din maximum 20.} \\
\texttt{Ați acumulat în mod colectiv 6 din 12 penalizări.} \\
\texttt{Celălalt jucător a trimis acest mesaj:} \\
\\ 
\\ 
\texttt{"Pe tabla mea, coloana 6 la rândurile 1{-}4 conține caractere de cadru: (6,1) e '│', (6,2) e '┘', (6,3) e '┼', (6,4) e '┌'. Nu pot muta acolo. Coloana 6 liberă e doar (6,6) și (6,7). Trebuie să găsim altă coloană. Cum arată coloana 7 la tine?"} \\
\\ 
\\ 
\texttt{Care este următoarea dvs. comandă?} \\
            }
        }
    }
    & \\ \\

    \theutterance \stepcounter{utterance}  
    & & & \multicolumn{4}{p{0.6\linewidth}}{
        \cellcolor[rgb]{0.9,0.9,0.9}{
            \makecell[{{p{\linewidth}}}]{
                \texttt{\tiny{[GM$\langle$P2]}}
                \texttt{SPUNE: Pe tabla mea, coloana 7: (7,1) ocupat de H, (7,2) e '◌' liber, (7,3) ocupat de C, (7,4) ocupat de T, (7,5) liber, (7,6) liber, (7,7) liber. Dar am caractere de cadru pe margine. Verific și coloana 1: (1,1) e '║' cadru, (1,2) e '╟' cadru, (1,3) e '║' cadru, (1,4) e '╟' cadru, (1,5) liber, (1,6) liber, (1,7) liber. Coloanele 2,3,4 la mine au multe cadre. Unde ai tu poziții libere? Poate ne întâlnim pe coloana 1 sau pe o coloană fără cadre pentru amândoi?} \\
            }
        }
    }
     \\ \\

    \theutterance \stepcounter{utterance}  
    & & & \multicolumn{2}{p{0.3\linewidth}}{
        \cellcolor[rgb]{0.9,0.9,0.9}{
            \makecell[{{p{\linewidth}}}]{
                \texttt{\tiny{[GM$|$GM]}}
                \texttt{SPUNE: Pe tabla mea, coloana 7: (7,1) ocupat de H, (7,2) e '◌' liber, (7,3) ocupat de C, (7,4) ocupat de T, (7,5) liber, (7,6) liber, (7,7) liber. Dar am caractere de cadru pe margine. Verific și coloana 1: (1,1) e '║' cadru, (1,2) e '╟' cadru, (1,3) e '║' cadru, (1,4) e '╟' cadru, (1,5) liber, (1,6) liber, (1,7) liber. Coloanele 2,3,4 la mine au multe cadre. Unde ai tu poziții libere? Poate ne întâlnim pe coloana 1 sau pe o coloană fără cadre pentru amândoi?} \\
            }
        }
    }
    & & \\ \\

    \theutterance \stepcounter{utterance}  
    & & \multicolumn{4}{p{0.6\linewidth}}{
        \cellcolor[rgb]{0.9,0.9,0.9}{
            \makecell[{{p{\linewidth}}}]{
                \texttt{\tiny{[P1$\langle$GM]}}
                \texttt{Mesajul dvs. a fost transmis celuilalt jucător.} \\
\\ 
\\ 
\texttt{În prezent jucați runda 6 din maximum 20.} \\
\texttt{Ați acumulat în mod colectiv 6 din 12 penalizări.} \\
\texttt{Celălalt jucător a trimis acest mesaj:} \\
\\ 
\\ 
\texttt{"Pe tabla mea, coloana 7: (7,1) ocupat de H, (7,2) e '◌' liber, (7,3) ocupat de C, (7,4) ocupat de T, (7,5) liber, (7,6) liber, (7,7) liber. Dar am caractere de cadru pe margine. Verific și coloana 1: (1,1) e '║' cadru, (1,2) e '╟' cadru, (1,3) e '║' cadru, (1,4) e '╟' cadru, (1,5) liber, (1,6) liber, (1,7) liber. Coloanele 2,3,4 la mine au multe cadre. Unde ai tu poziții libere? Poate ne întâlnim pe coloana 1 sau pe o coloană fără cadre pentru amândoi?"} \\
\\ 
\\ 
\texttt{Care este următoarea dvs. comandă?} \\
            }
        }
    }
    & \\ \\

    \theutterance \stepcounter{utterance}  
    & \multicolumn{4}{p{0.6\linewidth}}{
        \cellcolor[rgb]{0.9,0.9,0.9}{
            \makecell[{{p{\linewidth}}}]{
                \texttt{\tiny{[P1$\rangle$GM]}}
                \texttt{SPUNE: Pe tabla mea, coloana 1: (1,1) e '║', (1,2) e '║', (1,3) e '║', (1,4) e '║', (1,5) e '║', (1,6) ocupat de C, (1,7) e '║'. Coloana 2: (2,1) e '╤', (2,2) cadru, (2,3) cadru, (2,4) cadru, (2,5) liber, (2,6) liber, (2,7) ocupat de W. Coloana 3 are multe cadre. Coloana 4: (4,4) ocupat de H, restul cadre. Coloana 5: (5,7) ocupat de T, restul am caractere '│' și '┼'. Unde ai tu exact poziții libere fără cadre? Lista toate pozițiile goale de la tine.} \\
            }
        }
    }
    & & \\ \\

    \theutterance \stepcounter{utterance}  
    & & & \multicolumn{2}{p{0.3\linewidth}}{
        \cellcolor[rgb]{0.9,0.9,0.9}{
            \makecell[{{p{\linewidth}}}]{
                \texttt{\tiny{[GM$|$GM]}}
                \texttt{SPUNE: Pe tabla mea, coloana 1: (1,1) e '║', (1,2) e '║', (1,3) e '║', (1,4) e '║', (1,5) e '║', (1,6) ocupat de C, (1,7) e '║'. Coloana 2: (2,1) e '╤', (2,2) cadru, (2,3) cadru, (2,4) cadru, (2,5) liber, (2,6) liber, (2,7) ocupat de W. Coloana 3 are multe cadre. Coloana 4: (4,4) ocupat de H, restul cadre. Coloana 5: (5,7) ocupat de T, restul am caractere '│' și '┼'. Unde ai tu exact poziții libere fără cadre? Lista toate pozițiile goale de la tine.} \\
            }
        }
    }
    & & \\ \\

    \theutterance \stepcounter{utterance}  
    & & \multicolumn{4}{p{0.6\linewidth}}{
        \cellcolor[rgb]{0.9,0.9,0.9}{
            \makecell[{{p{\linewidth}}}]{
                \texttt{\tiny{[GM$\rangle$P2]}}
                \texttt{Mesajul dvs. a fost transmis celuilalt jucător.} \\
\\ 
\\ 
\texttt{În prezent jucați runda 7 din maximum 20.} \\
\texttt{Ați acumulat în mod colectiv 6 din 12 penalizări.} \\
\texttt{Celălalt jucător a trimis acest mesaj:} \\
\\ 
\\ 
\texttt{"Pe tabla mea, coloana 1: (1,1) e '║', (1,2) e '║', (1,3) e '║', (1,4) e '║', (1,5) e '║', (1,6) ocupat de C, (1,7) e '║'. Coloana 2: (2,1) e '╤', (2,2) cadru, (2,3) cadru, (2,4) cadru, (2,5) liber, (2,6) liber, (2,7) ocupat de W. Coloana 3 are multe cadre. Coloana 4: (4,4) ocupat de H, restul cadre. Coloana 5: (5,7) ocupat de T, restul am caractere '│' și '┼'. Unde ai tu exact poziții libere fără cadre? Lista toate pozițiile goale de la tine."} \\
\\ 
\\ 
\texttt{Care este următoarea dvs. comandă?} \\
            }
        }
    }
    & \\ \\

    \theutterance \stepcounter{utterance}  
    & & & \multicolumn{4}{p{0.6\linewidth}}{
        \cellcolor[rgb]{0.9,0.9,0.9}{
            \makecell[{{p{\linewidth}}}]{
                \texttt{\tiny{[GM$\langle$P2]}}
                \texttt{SPUNE: Pozițiile mele libere (cu '◌'): (2,1), (3,1), (4,1), (5,1), (6,1). (3,2), (4,2), (6,2). (2,3), (3,3), (4,3), (6,3). (2,4), (3,4), (6,4). (1,5) până la (7,5) toate libere. (1,6), (2,6), (3,6), (5,6), (6,6), (7,6) libere {-} (4,6) ocupat de I. (1,7) până la (4,7), (6,7), (7,7) libere {-} (5,7) ocupat de W. Rândul 5 e complet liber la mine! Putem pune toate obiectele pe rândul 5? W=(1,5), I=(2,5), T=(3,5), C=(4,5), H=(5,5)?} \\
            }
        }
    }
     \\ \\

    \theutterance \stepcounter{utterance}  
    & & & \multicolumn{2}{p{0.3\linewidth}}{
        \cellcolor[rgb]{0.9,0.9,0.9}{
            \makecell[{{p{\linewidth}}}]{
                \texttt{\tiny{[GM$|$GM]}}
                \texttt{SPUNE: Pozițiile mele libere (cu '◌'): (2,1), (3,1), (4,1), (5,1), (6,1). (3,2), (4,2), (6,2). (2,3), (3,3), (4,3), (6,3). (2,4), (3,4), (6,4). (1,5) până la (7,5) toate libere. (1,6), (2,6), (3,6), (5,6), (6,6), (7,6) libere {-} (4,6) ocupat de I. (1,7) până la (4,7), (6,7), (7,7) libere {-} (5,7) ocupat de W. Rândul 5 e complet liber la mine! Putem pune toate obiectele pe rândul 5? W=(1,5), I=(2,5), T=(3,5), C=(4,5), H=(5,5)?} \\
            }
        }
    }
    & & \\ \\

    \theutterance \stepcounter{utterance}  
    & & \multicolumn{4}{p{0.6\linewidth}}{
        \cellcolor[rgb]{0.9,0.9,0.9}{
            \makecell[{{p{\linewidth}}}]{
                \texttt{\tiny{[P1$\langle$GM]}}
                \texttt{Mesajul dvs. a fost transmis celuilalt jucător.} \\
\\ 
\\ 
\texttt{În prezent jucați runda 7 din maximum 20.} \\
\texttt{Ați acumulat în mod colectiv 6 din 12 penalizări.} \\
\texttt{Celălalt jucător a trimis acest mesaj:} \\
\\ 
\\ 
\texttt{"Pozițiile mele libere (cu '◌'): (2,1), (3,1), (4,1), (5,1), (6,1). (3,2), (4,2), (6,2). (2,3), (3,3), (4,3), (6,3). (2,4), (3,4), (6,4). (1,5) până la (7,5) toate libere. (1,6), (2,6), (3,6), (5,6), (6,6), (7,6) libere {-} (4,6) ocupat de I. (1,7) până la (4,7), (6,7), (7,7) libere {-} (5,7) ocupat de W. Rândul 5 e complet liber la mine! Putem pune toate obiectele pe rândul 5? W=(1,5), I=(2,5), T=(3,5), C=(4,5), H=(5,5)?"} \\
\\ 
\\ 
\texttt{Care este următoarea dvs. comandă?} \\
            }
        }
    }
    & \\ \\

    \theutterance \stepcounter{utterance}  
    & \multicolumn{4}{p{0.6\linewidth}}{
        \cellcolor[rgb]{0.9,0.9,0.9}{
            \makecell[{{p{\linewidth}}}]{
                \texttt{\tiny{[P1$\rangle$GM]}}
                \texttt{SPUNE: Pe tabla mea, rândul 5: (1,5) e '║', (2,5) e '◌' liber, (3,5) e '◌' liber, (4,5) e '◌' liber, (5,5) e '◌' liber, (6,5) ocupat de I, (7,5) e '║'. Deci (2,5), (3,5), (4,5), (5,5) sunt libere. Dar avem 5 obiecte și doar 4 poziții libere pe rândul 5. Putem folosi și rândul 6? La mine (2,6) liber, (3,6) liber, (4,6) liber, (5,6) liber, (6,6) liber, (7,6) liber. Propun: W=(2,5), I=(3,5), T=(4,5), C=(5,5), H=(2,6). Sau alt aranjament pe rândurile 5{-}6?} \\
            }
        }
    }
    & & \\ \\

    \theutterance \stepcounter{utterance}  
    & & & \multicolumn{2}{p{0.3\linewidth}}{
        \cellcolor[rgb]{0.9,0.9,0.9}{
            \makecell[{{p{\linewidth}}}]{
                \texttt{\tiny{[GM$|$GM]}}
                \texttt{SPUNE: Pe tabla mea, rândul 5: (1,5) e '║', (2,5) e '◌' liber, (3,5) e '◌' liber, (4,5) e '◌' liber, (5,5) e '◌' liber, (6,5) ocupat de I, (7,5) e '║'. Deci (2,5), (3,5), (4,5), (5,5) sunt libere. Dar avem 5 obiecte și doar 4 poziții libere pe rândul 5. Putem folosi și rândul 6? La mine (2,6) liber, (3,6) liber, (4,6) liber, (5,6) liber, (6,6) liber, (7,6) liber. Propun: W=(2,5), I=(3,5), T=(4,5), C=(5,5), H=(2,6). Sau alt aranjament pe rândurile 5{-}6?} \\
            }
        }
    }
    & & \\ \\

    \theutterance \stepcounter{utterance}  
    & & \multicolumn{4}{p{0.6\linewidth}}{
        \cellcolor[rgb]{0.9,0.9,0.9}{
            \makecell[{{p{\linewidth}}}]{
                \texttt{\tiny{[GM$\rangle$P2]}}
                \texttt{Mesajul dvs. a fost transmis celuilalt jucător.} \\
\\ 
\\ 
\texttt{În prezent jucați runda 8 din maximum 20.} \\
\texttt{Ați acumulat în mod colectiv 6 din 12 penalizări.} \\
\texttt{Celălalt jucător a trimis acest mesaj:} \\
\\ 
\\ 
\texttt{"Pe tabla mea, rândul 5: (1,5) e '║', (2,5) e '◌' liber, (3,5) e '◌' liber, (4,5) e '◌' liber, (5,5) e '◌' liber, (6,5) ocupat de I, (7,5) e '║'. Deci (2,5), (3,5), (4,5), (5,5) sunt libere. Dar avem 5 obiecte și doar 4 poziții libere pe rândul 5. Putem folosi și rândul 6? La mine (2,6) liber, (3,6) liber, (4,6) liber, (5,6) liber, (6,6) liber, (7,6) liber. Propun: W=(2,5), I=(3,5), T=(4,5), C=(5,5), H=(2,6). Sau alt aranjament pe rândurile 5{-}6?"} \\
\\ 
\\ 
\texttt{Care este următoarea dvs. comandă?} \\
            }
        }
    }
    & \\ \\

    \theutterance \stepcounter{utterance}  
    & & & \multicolumn{4}{p{0.6\linewidth}}{
        \cellcolor[rgb]{0.9,0.9,0.9}{
            \makecell[{{p{\linewidth}}}]{
                \texttt{\tiny{[GM$\langle$P2]}}
                \texttt{SPUNE: Bun! Pe tabla mea rândurile 5{-}6 sunt aproape toate libere. Accept: W=(2,5), I=(3,5), T=(4,5), C=(5,5), H=(2,6). Încep mutările acum!} \\
            }
        }
    }
     \\ \\

    \theutterance \stepcounter{utterance}  
    & & & \multicolumn{2}{p{0.3\linewidth}}{
        \cellcolor[rgb]{0.9,0.9,0.9}{
            \makecell[{{p{\linewidth}}}]{
                \texttt{\tiny{[GM$|$GM]}}
                \texttt{SPUNE: Bun! Pe tabla mea rândurile 5{-}6 sunt aproape toate libere. Accept: W=(2,5), I=(3,5), T=(4,5), C=(5,5), H=(2,6). Încep mutările acum!} \\
            }
        }
    }
    & & \\ \\

    \theutterance \stepcounter{utterance}  
    & & \multicolumn{4}{p{0.6\linewidth}}{
        \cellcolor[rgb]{0.9,0.9,0.9}{
            \makecell[{{p{\linewidth}}}]{
                \texttt{\tiny{[P1$\langle$GM]}}
                \texttt{Mesajul dvs. a fost transmis celuilalt jucător.} \\
\\ 
\\ 
\texttt{În prezent jucați runda 8 din maximum 20.} \\
\texttt{Ați acumulat în mod colectiv 6 din 12 penalizări.} \\
\texttt{Celălalt jucător a trimis acest mesaj:} \\
\\ 
\\ 
\texttt{"Bun! Pe tabla mea rândurile 5{-}6 sunt aproape toate libere. Accept: W=(2,5), I=(3,5), T=(4,5), C=(5,5), H=(2,6). Încep mutările acum!"} \\
\\ 
\\ 
\texttt{Care este următoarea dvs. comandă?} \\
            }
        }
    }
    & \\ \\

    \theutterance \stepcounter{utterance}  
    & \multicolumn{4}{p{0.6\linewidth}}{
        \cellcolor[rgb]{0.9,0.9,0.9}{
            \makecell[{{p{\linewidth}}}]{
                \texttt{\tiny{[P1$\rangle$GM]}}
                \texttt{MUTĂ: W, (2, 5)} \\
            }
        }
    }
    & & \\ \\

    \theutterance \stepcounter{utterance}  
    & & & \multicolumn{2}{p{0.3\linewidth}}{
        \cellcolor[rgb]{0.9,0.9,0.9}{
            \makecell[{{p{\linewidth}}}]{
                \texttt{\tiny{[GM$|$GM]}}
                \texttt{MUTĂ: W, (2, 5)} \\
            }
        }
    }
    & & \\ \\

    \theutterance \stepcounter{utterance}  
    & & & \multicolumn{2}{p{0.3\linewidth}}{
        \cellcolor[rgb]{0.9,0.9,0.9}{
            \makecell[{{p{\linewidth}}}]{
                \texttt{\tiny{[GM$|$GM]}}
                \texttt{W moved to (2, 5)} \\
            }
        }
    }
    & & \\ \\

    \theutterance \stepcounter{utterance}  
    & & \multicolumn{4}{p{0.6\linewidth}}{
        \cellcolor[rgb]{0.9,0.9,0.9}{
            \makecell[{{p{\linewidth}}}]{
                \texttt{\tiny{[GM$\rangle$P2]}}
                \texttt{Mesajul dvs. a fost transmis celuilalt jucător.} \\
\\ 
\\ 
\texttt{În prezent jucați runda 9 din maximum 20.} \\
\texttt{Ați acumulat în mod colectiv 6 din 12 penalizări.} \\
\texttt{Celălalt jucător a mutat un obiect pe grila sa.} \\
\\ 
\\ 
\texttt{Care este următoarea dvs. comandă?} \\
            }
        }
    }
    & \\ \\

    \theutterance \stepcounter{utterance}  
    & & & \multicolumn{4}{p{0.6\linewidth}}{
        \cellcolor[rgb]{0.9,0.9,0.9}{
            \makecell[{{p{\linewidth}}}]{
                \texttt{\tiny{[GM$\langle$P2]}}
                \texttt{MUTĂ: W, (2,5)} \\
            }
        }
    }
     \\ \\

    \theutterance \stepcounter{utterance}  
    & & & \multicolumn{2}{p{0.3\linewidth}}{
        \cellcolor[rgb]{0.9,0.9,0.9}{
            \makecell[{{p{\linewidth}}}]{
                \texttt{\tiny{[GM$|$GM]}}
                \texttt{MUTĂ: W, (2,5)} \\
            }
        }
    }
    & & \\ \\

    \theutterance \stepcounter{utterance}  
    & & & \multicolumn{2}{p{0.3\linewidth}}{
        \cellcolor[rgb]{0.9,0.9,0.9}{
            \makecell[{{p{\linewidth}}}]{
                \texttt{\tiny{[GM$|$GM]}}
                \texttt{W moved to (2, 5)} \\
            }
        }
    }
    & & \\ \\

    \theutterance \stepcounter{utterance}  
    & & \multicolumn{4}{p{0.6\linewidth}}{
        \cellcolor[rgb]{0.9,0.9,0.9}{
            \makecell[{{p{\linewidth}}}]{
                \texttt{\tiny{[P1$\langle$GM]}}
                \texttt{Obiectul 'W' a fost mutat la (2,5) cu succes. Grila dvs. actualizată arată astfel:} \\
\texttt{```} \\
\texttt{ 1234567} \\
\texttt{╔═╤═══╤═╗} \\
\texttt{║◌│◌◌◌│◌║ 1} \\
\texttt{╟─┼──┬┘◌║ 2} \\
\texttt{║◌│◌◌│◌◌║ 3} \\
\texttt{╟─┴──┘H◌║ 4} \\
\texttt{║◌W◌◌◌I◌║ 5} \\
\texttt{║C◌◌◌◌◌◌║ 6} \\
\texttt{║◌◌◌◌T◌◌║ 7} \\
\texttt{╚═══════╝} \\
\texttt{```} \\
\texttt{În prezent jucați runda 9 din maximum 20.} \\
\texttt{Ați acumulat în mod colectiv 6 din 12 penalizări.} \\
\texttt{Celălalt jucător a mutat un obiect pe grila sa.} \\
\\ 
\\ 
\texttt{Care este următoarea dvs. comandă?} \\
            }
        }
    }
    & \\ \\

    \theutterance \stepcounter{utterance}  
    & \multicolumn{4}{p{0.6\linewidth}}{
        \cellcolor[rgb]{0.9,0.9,0.9}{
            \makecell[{{p{\linewidth}}}]{
                \texttt{\tiny{[P1$\rangle$GM]}}
                \texttt{MUTĂ: I, (3, 5)} \\
            }
        }
    }
    & & \\ \\

    \theutterance \stepcounter{utterance}  
    & & & \multicolumn{2}{p{0.3\linewidth}}{
        \cellcolor[rgb]{0.9,0.9,0.9}{
            \makecell[{{p{\linewidth}}}]{
                \texttt{\tiny{[GM$|$GM]}}
                \texttt{MUTĂ: I, (3, 5)} \\
            }
        }
    }
    & & \\ \\

    \theutterance \stepcounter{utterance}  
    & & & \multicolumn{2}{p{0.3\linewidth}}{
        \cellcolor[rgb]{0.9,0.9,0.9}{
            \makecell[{{p{\linewidth}}}]{
                \texttt{\tiny{[GM$|$GM]}}
                \texttt{I moved to (3, 5)} \\
            }
        }
    }
    & & \\ \\

    \theutterance \stepcounter{utterance}  
    & & \multicolumn{4}{p{0.6\linewidth}}{
        \cellcolor[rgb]{0.9,0.9,0.9}{
            \makecell[{{p{\linewidth}}}]{
                \texttt{\tiny{[GM$\rangle$P2]}}
                \texttt{Obiectul 'W' a fost mutat la (2,5) cu succes. Grila dvs. actualizată arată astfel:} \\
\texttt{```} \\
\texttt{ 1234567} \\
\texttt{╔═╤═══╤═╗} \\
\texttt{║◌│◌◌◌│H║ 1} \\
\texttt{╟─┼──┬┘◌║ 2} \\
\texttt{║◌│◌◌│◌C║ 3} \\
\texttt{╟─┴──┘◌T║ 4} \\
\texttt{║◌W◌◌◌◌◌║ 5} \\
\texttt{║◌◌◌I◌◌◌║ 6} \\
\texttt{║◌◌◌◌◌◌◌║ 7} \\
\texttt{╚═══════╝} \\
\texttt{```} \\
\texttt{În prezent jucați runda 10 din maximum 20.} \\
\texttt{Ați acumulat în mod colectiv 6 din 12 penalizări.} \\
\texttt{Celălalt jucător a mutat un obiect pe grila sa.} \\
\\ 
\\ 
\texttt{Care este următoarea dvs. comandă?} \\
            }
        }
    }
    & \\ \\

    \theutterance \stepcounter{utterance}  
    & & & \multicolumn{4}{p{0.6\linewidth}}{
        \cellcolor[rgb]{0.9,0.9,0.9}{
            \makecell[{{p{\linewidth}}}]{
                \texttt{\tiny{[GM$\langle$P2]}}
                \texttt{MUTĂ: I, (3,5)} \\
            }
        }
    }
     \\ \\

    \theutterance \stepcounter{utterance}  
    & & & \multicolumn{2}{p{0.3\linewidth}}{
        \cellcolor[rgb]{0.9,0.9,0.9}{
            \makecell[{{p{\linewidth}}}]{
                \texttt{\tiny{[GM$|$GM]}}
                \texttt{MUTĂ: I, (3,5)} \\
            }
        }
    }
    & & \\ \\

    \theutterance \stepcounter{utterance}  
    & & & \multicolumn{2}{p{0.3\linewidth}}{
        \cellcolor[rgb]{0.9,0.9,0.9}{
            \makecell[{{p{\linewidth}}}]{
                \texttt{\tiny{[GM$|$GM]}}
                \texttt{I moved to (3, 5)} \\
            }
        }
    }
    & & \\ \\

    \theutterance \stepcounter{utterance}  
    & & \multicolumn{4}{p{0.6\linewidth}}{
        \cellcolor[rgb]{0.9,0.9,0.9}{
            \makecell[{{p{\linewidth}}}]{
                \texttt{\tiny{[P1$\langle$GM]}}
                \texttt{Obiectul 'I' a fost mutat la (3,5) cu succes. Grila dvs. actualizată arată astfel:} \\
\texttt{```} \\
\texttt{ 1234567} \\
\texttt{╔═╤═══╤═╗} \\
\texttt{║◌│◌◌◌│◌║ 1} \\
\texttt{╟─┼──┬┘◌║ 2} \\
\texttt{║◌│◌◌│◌◌║ 3} \\
\texttt{╟─┴──┘H◌║ 4} \\
\texttt{║◌WI◌◌◌◌║ 5} \\
\texttt{║C◌◌◌◌◌◌║ 6} \\
\texttt{║◌◌◌◌T◌◌║ 7} \\
\texttt{╚═══════╝} \\
\texttt{```} \\
\texttt{În prezent jucați runda 10 din maximum 20.} \\
\texttt{Ați acumulat în mod colectiv 6 din 12 penalizări.} \\
\texttt{Celălalt jucător a mutat un obiect pe grila sa.} \\
\\ 
\\ 
\texttt{Care este următoarea dvs. comandă?} \\
            }
        }
    }
    & \\ \\

    \theutterance \stepcounter{utterance}  
    & \multicolumn{4}{p{0.6\linewidth}}{
        \cellcolor[rgb]{0.9,0.9,0.9}{
            \makecell[{{p{\linewidth}}}]{
                \texttt{\tiny{[P1$\rangle$GM]}}
                \texttt{MUTĂ: T, (4, 5)} \\
            }
        }
    }
    & & \\ \\

    \theutterance \stepcounter{utterance}  
    & & & \multicolumn{2}{p{0.3\linewidth}}{
        \cellcolor[rgb]{0.9,0.9,0.9}{
            \makecell[{{p{\linewidth}}}]{
                \texttt{\tiny{[GM$|$GM]}}
                \texttt{MUTĂ: T, (4, 5)} \\
            }
        }
    }
    & & \\ \\

    \theutterance \stepcounter{utterance}  
    & & & \multicolumn{2}{p{0.3\linewidth}}{
        \cellcolor[rgb]{0.9,0.9,0.9}{
            \makecell[{{p{\linewidth}}}]{
                \texttt{\tiny{[GM$|$GM]}}
                \texttt{T moved to (4, 5)} \\
            }
        }
    }
    & & \\ \\

    \theutterance \stepcounter{utterance}  
    & & \multicolumn{4}{p{0.6\linewidth}}{
        \cellcolor[rgb]{0.9,0.9,0.9}{
            \makecell[{{p{\linewidth}}}]{
                \texttt{\tiny{[GM$\rangle$P2]}}
                \texttt{Obiectul 'I' a fost mutat la (3,5) cu succes. Grila dvs. actualizată arată astfel:} \\
\texttt{```} \\
\texttt{ 1234567} \\
\texttt{╔═╤═══╤═╗} \\
\texttt{║◌│◌◌◌│H║ 1} \\
\texttt{╟─┼──┬┘◌║ 2} \\
\texttt{║◌│◌◌│◌C║ 3} \\
\texttt{╟─┴──┘◌T║ 4} \\
\texttt{║◌WI◌◌◌◌║ 5} \\
\texttt{║◌◌◌◌◌◌◌║ 6} \\
\texttt{║◌◌◌◌◌◌◌║ 7} \\
\texttt{╚═══════╝} \\
\texttt{```} \\
\texttt{În prezent jucați runda 11 din maximum 20.} \\
\texttt{Ați acumulat în mod colectiv 6 din 12 penalizări.} \\
\texttt{Celălalt jucător a mutat un obiect pe grila sa.} \\
\\ 
\\ 
\texttt{Care este următoarea dvs. comandă?} \\
            }
        }
    }
    & \\ \\

    \theutterance \stepcounter{utterance}  
    & & & \multicolumn{4}{p{0.6\linewidth}}{
        \cellcolor[rgb]{0.9,0.9,0.9}{
            \makecell[{{p{\linewidth}}}]{
                \texttt{\tiny{[GM$\langle$P2]}}
                \texttt{MUTĂ: T, (4,5)} \\
            }
        }
    }
     \\ \\

    \theutterance \stepcounter{utterance}  
    & & & \multicolumn{2}{p{0.3\linewidth}}{
        \cellcolor[rgb]{0.9,0.9,0.9}{
            \makecell[{{p{\linewidth}}}]{
                \texttt{\tiny{[GM$|$GM]}}
                \texttt{MUTĂ: T, (4,5)} \\
            }
        }
    }
    & & \\ \\

    \theutterance \stepcounter{utterance}  
    & & & \multicolumn{2}{p{0.3\linewidth}}{
        \cellcolor[rgb]{0.9,0.9,0.9}{
            \makecell[{{p{\linewidth}}}]{
                \texttt{\tiny{[GM$|$GM]}}
                \texttt{T moved to (4, 5)} \\
            }
        }
    }
    & & \\ \\

    \theutterance \stepcounter{utterance}  
    & & \multicolumn{4}{p{0.6\linewidth}}{
        \cellcolor[rgb]{0.9,0.9,0.9}{
            \makecell[{{p{\linewidth}}}]{
                \texttt{\tiny{[P1$\langle$GM]}}
                \texttt{Obiectul 'T' a fost mutat la (4,5) cu succes. Grila dvs. actualizată arată astfel:} \\
\texttt{```} \\
\texttt{ 1234567} \\
\texttt{╔═╤═══╤═╗} \\
\texttt{║◌│◌◌◌│◌║ 1} \\
\texttt{╟─┼──┬┘◌║ 2} \\
\texttt{║◌│◌◌│◌◌║ 3} \\
\texttt{╟─┴──┘H◌║ 4} \\
\texttt{║◌WIT◌◌◌║ 5} \\
\texttt{║C◌◌◌◌◌◌║ 6} \\
\texttt{║◌◌◌◌◌◌◌║ 7} \\
\texttt{╚═══════╝} \\
\texttt{```} \\
\texttt{În prezent jucați runda 11 din maximum 20.} \\
\texttt{Ați acumulat în mod colectiv 6 din 12 penalizări.} \\
\texttt{Celălalt jucător a mutat un obiect pe grila sa.} \\
\\ 
\\ 
\texttt{Care este următoarea dvs. comandă?} \\
            }
        }
    }
    & \\ \\

    \theutterance \stepcounter{utterance}  
    & \multicolumn{4}{p{0.6\linewidth}}{
        \cellcolor[rgb]{0.9,0.9,0.9}{
            \makecell[{{p{\linewidth}}}]{
                \texttt{\tiny{[P1$\rangle$GM]}}
                \texttt{MUTĂ: C, (5, 5)} \\
            }
        }
    }
    & & \\ \\

    \theutterance \stepcounter{utterance}  
    & & & \multicolumn{2}{p{0.3\linewidth}}{
        \cellcolor[rgb]{0.9,0.9,0.9}{
            \makecell[{{p{\linewidth}}}]{
                \texttt{\tiny{[GM$|$GM]}}
                \texttt{MUTĂ: C, (5, 5)} \\
            }
        }
    }
    & & \\ \\

    \theutterance \stepcounter{utterance}  
    & & & \multicolumn{2}{p{0.3\linewidth}}{
        \cellcolor[rgb]{0.9,0.9,0.9}{
            \makecell[{{p{\linewidth}}}]{
                \texttt{\tiny{[GM$|$GM]}}
                \texttt{C moved to (5, 5)} \\
            }
        }
    }
    & & \\ \\

    \theutterance \stepcounter{utterance}  
    & & \multicolumn{4}{p{0.6\linewidth}}{
        \cellcolor[rgb]{0.9,0.9,0.9}{
            \makecell[{{p{\linewidth}}}]{
                \texttt{\tiny{[GM$\rangle$P2]}}
                \texttt{Obiectul 'T' a fost mutat la (4,5) cu succes. Grila dvs. actualizată arată astfel:} \\
\texttt{```} \\
\texttt{ 1234567} \\
\texttt{╔═╤═══╤═╗} \\
\texttt{║◌│◌◌◌│H║ 1} \\
\texttt{╟─┼──┬┘◌║ 2} \\
\texttt{║◌│◌◌│◌C║ 3} \\
\texttt{╟─┴──┘◌◌║ 4} \\
\texttt{║◌WIT◌◌◌║ 5} \\
\texttt{║◌◌◌◌◌◌◌║ 6} \\
\texttt{║◌◌◌◌◌◌◌║ 7} \\
\texttt{╚═══════╝} \\
\texttt{```} \\
\texttt{În prezent jucați runda 12 din maximum 20.} \\
\texttt{Ați acumulat în mod colectiv 6 din 12 penalizări.} \\
\texttt{Celălalt jucător a mutat un obiect pe grila sa.} \\
\\ 
\\ 
\texttt{Care este următoarea dvs. comandă?} \\
            }
        }
    }
    & \\ \\

    \theutterance \stepcounter{utterance}  
    & & & \multicolumn{4}{p{0.6\linewidth}}{
        \cellcolor[rgb]{0.9,0.9,0.9}{
            \makecell[{{p{\linewidth}}}]{
                \texttt{\tiny{[GM$\langle$P2]}}
                \texttt{MUTĂ: C, (5,5)} \\
            }
        }
    }
     \\ \\

    \theutterance \stepcounter{utterance}  
    & & & \multicolumn{2}{p{0.3\linewidth}}{
        \cellcolor[rgb]{0.9,0.9,0.9}{
            \makecell[{{p{\linewidth}}}]{
                \texttt{\tiny{[GM$|$GM]}}
                \texttt{MUTĂ: C, (5,5)} \\
            }
        }
    }
    & & \\ \\

    \theutterance \stepcounter{utterance}  
    & & & \multicolumn{2}{p{0.3\linewidth}}{
        \cellcolor[rgb]{0.9,0.9,0.9}{
            \makecell[{{p{\linewidth}}}]{
                \texttt{\tiny{[GM$|$GM]}}
                \texttt{C moved to (5, 5)} \\
            }
        }
    }
    & & \\ \\

    \theutterance \stepcounter{utterance}  
    & & \multicolumn{4}{p{0.6\linewidth}}{
        \cellcolor[rgb]{0.9,0.9,0.9}{
            \makecell[{{p{\linewidth}}}]{
                \texttt{\tiny{[P1$\langle$GM]}}
                \texttt{Obiectul 'C' a fost mutat la (5,5) cu succes. Grila dvs. actualizată arată astfel:} \\
\texttt{```} \\
\texttt{ 1234567} \\
\texttt{╔═╤═══╤═╗} \\
\texttt{║◌│◌◌◌│◌║ 1} \\
\texttt{╟─┼──┬┘◌║ 2} \\
\texttt{║◌│◌◌│◌◌║ 3} \\
\texttt{╟─┴──┘H◌║ 4} \\
\texttt{║◌WITC◌◌║ 5} \\
\texttt{║◌◌◌◌◌◌◌║ 6} \\
\texttt{║◌◌◌◌◌◌◌║ 7} \\
\texttt{╚═══════╝} \\
\texttt{```} \\
\texttt{În prezent jucați runda 12 din maximum 20.} \\
\texttt{Ați acumulat în mod colectiv 6 din 12 penalizări.} \\
\texttt{Celălalt jucător a mutat un obiect pe grila sa.} \\
\\ 
\\ 
\texttt{Care este următoarea dvs. comandă?} \\
            }
        }
    }
    & \\ \\

    \theutterance \stepcounter{utterance}  
    & \multicolumn{4}{p{0.6\linewidth}}{
        \cellcolor[rgb]{0.9,0.9,0.9}{
            \makecell[{{p{\linewidth}}}]{
                \texttt{\tiny{[P1$\rangle$GM]}}
                \texttt{MUTĂ: H, (2, 6)} \\
            }
        }
    }
    & & \\ \\

    \theutterance \stepcounter{utterance}  
    & & & \multicolumn{2}{p{0.3\linewidth}}{
        \cellcolor[rgb]{0.9,0.9,0.9}{
            \makecell[{{p{\linewidth}}}]{
                \texttt{\tiny{[GM$|$GM]}}
                \texttt{MUTĂ: H, (2, 6)} \\
            }
        }
    }
    & & \\ \\

    \theutterance \stepcounter{utterance}  
    & & & \multicolumn{2}{p{0.3\linewidth}}{
        \cellcolor[rgb]{0.9,0.9,0.9}{
            \makecell[{{p{\linewidth}}}]{
                \texttt{\tiny{[GM$|$GM]}}
                \texttt{H moved to (2, 6)} \\
            }
        }
    }
    & & \\ \\

    \theutterance \stepcounter{utterance}  
    & & \multicolumn{4}{p{0.6\linewidth}}{
        \cellcolor[rgb]{0.9,0.9,0.9}{
            \makecell[{{p{\linewidth}}}]{
                \texttt{\tiny{[GM$\rangle$P2]}}
                \texttt{Obiectul 'C' a fost mutat la (5,5) cu succes. Grila dvs. actualizată arată astfel:} \\
\texttt{```} \\
\texttt{ 1234567} \\
\texttt{╔═╤═══╤═╗} \\
\texttt{║◌│◌◌◌│H║ 1} \\
\texttt{╟─┼──┬┘◌║ 2} \\
\texttt{║◌│◌◌│◌◌║ 3} \\
\texttt{╟─┴──┘◌◌║ 4} \\
\texttt{║◌WITC◌◌║ 5} \\
\texttt{║◌◌◌◌◌◌◌║ 6} \\
\texttt{║◌◌◌◌◌◌◌║ 7} \\
\texttt{╚═══════╝} \\
\texttt{```} \\
\texttt{În prezent jucați runda 13 din maximum 20.} \\
\texttt{Ați acumulat în mod colectiv 6 din 12 penalizări.} \\
\texttt{Celălalt jucător a mutat un obiect pe grila sa.} \\
\\ 
\\ 
\texttt{Care este următoarea dvs. comandă?} \\
            }
        }
    }
    & \\ \\

    \theutterance \stepcounter{utterance}  
    & & & \multicolumn{4}{p{0.6\linewidth}}{
        \cellcolor[rgb]{0.9,0.9,0.9}{
            \makecell[{{p{\linewidth}}}]{
                \texttt{\tiny{[GM$\langle$P2]}}
                \texttt{MUTĂ: H, (2,6)} \\
            }
        }
    }
     \\ \\

    \theutterance \stepcounter{utterance}  
    & & & \multicolumn{2}{p{0.3\linewidth}}{
        \cellcolor[rgb]{0.9,0.9,0.9}{
            \makecell[{{p{\linewidth}}}]{
                \texttt{\tiny{[GM$|$GM]}}
                \texttt{MUTĂ: H, (2,6)} \\
            }
        }
    }
    & & \\ \\

    \theutterance \stepcounter{utterance}  
    & & & \multicolumn{2}{p{0.3\linewidth}}{
        \cellcolor[rgb]{0.9,0.9,0.9}{
            \makecell[{{p{\linewidth}}}]{
                \texttt{\tiny{[GM$|$GM]}}
                \texttt{H moved to (2, 6)} \\
            }
        }
    }
    & & \\ \\

    \theutterance \stepcounter{utterance}  
    & & \multicolumn{4}{p{0.6\linewidth}}{
        \cellcolor[rgb]{0.9,0.9,0.9}{
            \makecell[{{p{\linewidth}}}]{
                \texttt{\tiny{[P1$\langle$GM]}}
                \texttt{Obiectul 'H' a fost mutat la (2,6) cu succes. Grila dvs. actualizată arată astfel:} \\
\texttt{```} \\
\texttt{ 1234567} \\
\texttt{╔═╤═══╤═╗} \\
\texttt{║◌│◌◌◌│◌║ 1} \\
\texttt{╟─┼──┬┘◌║ 2} \\
\texttt{║◌│◌◌│◌◌║ 3} \\
\texttt{╟─┴──┘◌◌║ 4} \\
\texttt{║◌WITC◌◌║ 5} \\
\texttt{║◌H◌◌◌◌◌║ 6} \\
\texttt{║◌◌◌◌◌◌◌║ 7} \\
\texttt{╚═══════╝} \\
\texttt{```} \\
\texttt{În prezent jucați runda 13 din maximum 20.} \\
\texttt{Ați acumulat în mod colectiv 6 din 12 penalizări.} \\
\texttt{Celălalt jucător a mutat un obiect pe grila sa.} \\
\\ 
\\ 
\texttt{Care este următoarea dvs. comandă?} \\
            }
        }
    }
    & \\ \\

    \theutterance \stepcounter{utterance}  
    & \multicolumn{4}{p{0.6\linewidth}}{
        \cellcolor[rgb]{0.9,0.9,0.9}{
            \makecell[{{p{\linewidth}}}]{
                \texttt{\tiny{[P1$\rangle$GM]}}
                \texttt{SPUNE: Am mutat toate obiectele pe pozițiile convenite: W=(2,5), I=(3,5), T=(4,5), C=(5,5), H=(2,6). Tu ai terminat mutările? terminat?} \\
            }
        }
    }
    & & \\ \\

    \theutterance \stepcounter{utterance}  
    & & & \multicolumn{2}{p{0.3\linewidth}}{
        \cellcolor[rgb]{0.9,0.9,0.9}{
            \makecell[{{p{\linewidth}}}]{
                \texttt{\tiny{[GM$|$GM]}}
                \texttt{SPUNE: Am mutat toate obiectele pe pozițiile convenite: W=(2,5), I=(3,5), T=(4,5), C=(5,5), H=(2,6). Tu ai terminat mutările? terminat?} \\
            }
        }
    }
    & & \\ \\

    \theutterance \stepcounter{utterance}  
    & & \multicolumn{4}{p{0.6\linewidth}}{
        \cellcolor[rgb]{0.9,0.9,0.9}{
            \makecell[{{p{\linewidth}}}]{
                \texttt{\tiny{[GM$\rangle$P2]}}
                \texttt{Obiectul 'H' a fost mutat la (2,6) cu succes. Grila dvs. actualizată arată astfel:} \\
\texttt{```} \\
\texttt{ 1234567} \\
\texttt{╔═╤═══╤═╗} \\
\texttt{║◌│◌◌◌│◌║ 1} \\
\texttt{╟─┼──┬┘◌║ 2} \\
\texttt{║◌│◌◌│◌◌║ 3} \\
\texttt{╟─┴──┘◌◌║ 4} \\
\texttt{║◌WITC◌◌║ 5} \\
\texttt{║◌H◌◌◌◌◌║ 6} \\
\texttt{║◌◌◌◌◌◌◌║ 7} \\
\texttt{╚═══════╝} \\
\texttt{```} \\
\texttt{În prezent jucați runda 14 din maximum 20.} \\
\texttt{Ați acumulat în mod colectiv 6 din 12 penalizări.} \\
\texttt{Celălalt jucător a trimis acest mesaj:} \\
\\ 
\\ 
\texttt{"Am mutat toate obiectele pe pozițiile convenite: W=(2,5), I=(3,5), T=(4,5), C=(5,5), H=(2,6). Tu ai terminat mutările? terminat?"} \\
\\ 
\\ 
\texttt{Care este următoarea dvs. comandă?} \\
            }
        }
    }
    & \\ \\

    \theutterance \stepcounter{utterance}  
    & & & \multicolumn{4}{p{0.6\linewidth}}{
        \cellcolor[rgb]{0.9,0.9,0.9}{
            \makecell[{{p{\linewidth}}}]{
                \texttt{\tiny{[GM$\langle$P2]}}
                \texttt{SPUNE: Și eu am terminat toate mutările! W=(2,5), I=(3,5), T=(4,5), C=(5,5), H=(2,6). terminat!} \\
            }
        }
    }
     \\ \\

    \theutterance \stepcounter{utterance}  
    & & & \multicolumn{2}{p{0.3\linewidth}}{
        \cellcolor[rgb]{0.9,0.9,0.9}{
            \makecell[{{p{\linewidth}}}]{
                \texttt{\tiny{[GM$|$GM]}}
                \texttt{SPUNE: Și eu am terminat toate mutările! W=(2,5), I=(3,5), T=(4,5), C=(5,5), H=(2,6). terminat!} \\
            }
        }
    }
    & & \\ \\

    \theutterance \stepcounter{utterance}  
    & & & \multicolumn{2}{p{0.3\linewidth}}{
        \cellcolor[rgb]{0.9,0.9,0.9}{
            \makecell[{{p{\linewidth}}}]{
                \texttt{\tiny{[GM$|$GM]}}
                \texttt{true} \\
            }
        }
    }
    & & \\ \\

    \theutterance \stepcounter{utterance}  
    & & & \multicolumn{2}{p{0.3\linewidth}}{
        \cellcolor[rgb]{0.9,0.9,0.9}{
            \makecell[{{p{\linewidth}}}]{
                \texttt{\tiny{[GM$|$GM]}}
                \texttt{Initial states:} \\
\texttt{```} \\
\texttt{Player 1:} \\
\texttt{ 1234567} \\
\texttt{╔═╤═══╤═╗} \\
\texttt{║◌│◌◌◌│◌║ 1} \\
\texttt{╟─┼──┬┘◌║ 2} \\
\texttt{║◌│◌◌│◌◌║ 3} \\
\texttt{╟─┴──┘H◌║ 4} \\
\texttt{║◌◌◌◌◌I◌║ 5} \\
\texttt{║C◌◌◌◌◌◌║ 6} \\
\texttt{║◌W◌◌T◌◌║ 7} \\
\texttt{╚═══════╝} \\
\\ 
\\ 
\texttt{Player 2:} \\
\texttt{ 1234567} \\
\texttt{╔═╤═══╤═╗} \\
\texttt{║◌│◌◌◌│H║ 1} \\
\texttt{╟─┼──┬┘◌║ 2} \\
\texttt{║◌│◌◌│◌C║ 3} \\
\texttt{╟─┴──┘◌T║ 4} \\
\texttt{║◌◌◌◌◌◌◌║ 5} \\
\texttt{║◌◌◌I◌◌◌║ 6} \\
\texttt{║◌◌◌◌W◌◌║ 7} \\
\texttt{╚═══════╝} \\
\\ 
\texttt{```} \\
            }
        }
    }
    & & \\ \\

    \theutterance \stepcounter{utterance}  
    & & & \multicolumn{2}{p{0.3\linewidth}}{
        \cellcolor[rgb]{0.9,0.9,0.9}{
            \makecell[{{p{\linewidth}}}]{
                \texttt{\tiny{[GM$|$GM]}}
                \texttt{End states:} \\
\texttt{```} \\
\texttt{Player 1:} \\
\texttt{ 1234567} \\
\texttt{╔═╤═══╤═╗} \\
\texttt{║◌│◌◌◌│◌║ 1} \\
\texttt{╟─┼──┬┘◌║ 2} \\
\texttt{║◌│◌◌│◌◌║ 3} \\
\texttt{╟─┴──┘◌◌║ 4} \\
\texttt{║◌WITC◌◌║ 5} \\
\texttt{║◌H◌◌◌◌◌║ 6} \\
\texttt{║◌◌◌◌◌◌◌║ 7} \\
\texttt{╚═══════╝} \\
\\ 
\\ 
\texttt{Player 2:} \\
\texttt{ 1234567} \\
\texttt{╔═╤═══╤═╗} \\
\texttt{║◌│◌◌◌│◌║ 1} \\
\texttt{╟─┼──┬┘◌║ 2} \\
\texttt{║◌│◌◌│◌◌║ 3} \\
\texttt{╟─┴──┘◌◌║ 4} \\
\texttt{║◌WITC◌◌║ 5} \\
\texttt{║◌H◌◌◌◌◌║ 6} \\
\texttt{║◌◌◌◌◌◌◌║ 7} \\
\texttt{╚═══════╝} \\
\\ 
\texttt{```} \\
            }
        }
    }
    & & \\ \\

    \theutterance \stepcounter{utterance}  
    & & & \multicolumn{2}{p{0.3\linewidth}}{
        \cellcolor[rgb]{0.9,0.9,0.9}{
            \makecell[{{p{\linewidth}}}]{
                \texttt{\tiny{[GM$|$GM]}}
                \texttt{* success: True} \\
\texttt{* lose: False} \\
\texttt{* aborted: False} \\
\texttt{{-}{-}{-}{-}{-}{-}{-}} \\
\texttt{* Shifts: 4.00} \\
\texttt{* Max Shifts: 8.00} \\
\texttt{* Min Shifts: 4.00} \\
\texttt{* End Distance Sum: 0.00} \\
\texttt{* Init Distance Sum: 18.71} \\
\texttt{* Expected Distance Sum: 20.95} \\
\texttt{* Penalties: 6.00} \\
\texttt{* Max Penalties: 12.00} \\
\texttt{* Rounds: 14.00} \\
\texttt{* Max Rounds: 20.00} \\
\texttt{* Object Count: 5.00} \\
            }
        }
    }
    & & \\ \\

\end{supertabular}
}

\end{document}
